\subsection{Case III}
\label{sec:molecular-algebraic-induction-caseIII}

Case III realizes the reduction move left by Cases~I and~II: splitting off a
degree-$2$ vertex $v$ of a simple minimal $0$-dof-graph $G$ that has no proper
rigid subgraph (KT Lemmas~6.10/6.13, \cref{lem:case-III}). The placement of
\cref{lem:case-II-realization-placement} realizes $G$ at rank $D(|V| - 1) - 1$,
one row short of full rank, and Case III supplies the missing row. That row is
KT Claim~6.11 (eq.\ (6.23)): in the split-off graph $G_v^{ab}$, one rigidity row
at the deleted edge $ab$ is a linear combination of the others, so it is
redundant. This subsection proves Claim~6.11; the candidate-completion
then rewrites the redundant row into the missing full-rank row, and
Claim~6.12 discharges the one remaining conditional.

\subsubsection{KT Claim~6.11: a redundant row at the new edge}

In Katoh and Tanigawa's proof, the induction hypothesis gives a generic
realization $(G_v^{ab}, q)$ of $G_v^{ab}$ of rank $D(|V(G)| - 2)$ (KT
eq.~(6.18), p.~682). They call a row of $R(G_v^{ab}, q)$ redundant if
removing it does not lower the rank. KT Claim~6.11 says that one of the
$D - 1$ rows at the new edge $ab$ is
redundant~\cite[p.~684]{katohTanigawa2011}. Its proof finds the
redundancy first in the matroid $M(\widetilde{G_v^{ab}})$ and then in
the rows, in three steps. First, $M(\widetilde{G_v^{ab}})$ has a base
that does not contain every copy of $ab$
(\cref{lem:case-III-claim-6-11-base}). Second, it follows that the
graph $G_v = G - v$, which is $G_v^{ab}$ without the edge $ab$, is a
minimal $k'$-dof-graph with $k' \le D - 2$
(\cref{lem:case-III-gap3-minimalKDof}). So the induction hypothesis
gives $\operatorname{rank} R(G_v, q) \ge D(|V(G)| - 2) - (D - 2)$ (KT
eq.~(6.22), \cref{lem:case-III-nested-rank-lower}). Third, the rows of
$R(G_v^{ab}, q)$ are those of $R(G_v, q)$ and the rows at $ab$. The rows
at $ab$ therefore raise the rank by at most $D - 2$, and one of them is
redundant (KT eq.~(6.23), \cref{lem:case-III-claim-6-11}).

The first step is the forward half of KT Lemma~4.3(ii), which
\cref{lem:splitoff-kdof-criterion} states with its converse. We state it
on its own, as Katoh and Tanigawa use it here.

\begin{lemma}[A base of $M(\widetilde{G_v^{ab}})$ without some copy of $ab$; KT Lemma~4.3(ii)]
  \label{lem:case-III-claim-6-11-base}
  \lean{Graph.splitOff_exists_base_inter_fiber_lt}
  \leanok
  \uses{def:k-dof, def:matroid-MG, thm:def-eq-corank, lem:forest-surgery-count, lem:splitoff-deficiency}
  Let $D \ge 2$, and let $G$ be a $k$-dof-graph with a vertex $v$ of
  degree two, whose edges $e_a$ and $e_b$ join $v$ to distinct vertices
  $a, b \neq v$. Let $\widetilde{ab}$ be the set of the $D - 1$ copies
  in $\widetilde{G_v^{ab}}$ of the new edge $ab$ of $G_v^{ab}$. If
  $G_v^{ab}$ is a $k$-dof-graph, then $M(\widetilde{G_v^{ab}})$ has a
  base $B'$ with
  \[
    \bigl|\widetilde{ab} \cap B'\bigr| \;<\; D - 1.
  \]
\end{lemma}
\begin{proof}
  \leanok
  We recall Katoh and Tanigawa's
  proof~\cite[p.~661]{katohTanigawa2011}, which
  \cref{lem:splitoff-kdof-criterion} also gives. Let $B$ be a base of
  $M(\tilde G)$. By \cref{lem:balanced-forest-packing}, $B$ is the union
  of $D$ pairwise disjoint forests of $\tilde G$, each of which has an
  edge incident with $v$, as the proof of KT Lemma~4.1 assumes
  (\cref{rem:kt-lemma-41}). The surgery of KT Lemma~4.1 on these forests
  (\cref{lem:forest-surgery-count}) gives an independent set $I'$ of
  $M(\widetilde{G_v^{ab}})$ with $|I'| = |B| - D$ and
  $|\widetilde{ab} \cap I'| < D - 1$.

  The vertices of $G_v^{ab}$ are the $|V(G)| - 1$ vertices of $G$ other
  than $v$. Both graphs are $k$-dof-graphs, so KT eq.~(2.4)
  (\cref{thm:def-eq-corank}) gives $|B| = D(|V(G)| - 1) - k$ and
  \[
    \operatorname{rank} M(\widetilde{G_v^{ab}})
    \;=\; D\,(|V(G)| - 2) - k \;=\; |B| - D \;=\; |I'|.
  \]
  So $I'$ is a base, and we take $B' = I'$. A copy of $ab$ outside $B'$
  can be removed from $\widetilde{G_v^{ab}}$ without lowering the rank of
  its matroid. Katoh and Tanigawa call such a copy a redundant edge
  (p.~684).
\end{proof}

\begin{fmlnote}
  The pinned theorem is also the first pinned theorem of
  \cref{lem:splitoff-kdof-criterion}. It assumes that the label of the
  new edge $ab$ lies outside $E(G)$ (\cref{fmlnote:graph-operations}).
\end{fmlnote}

In the second step, $G$ is a minimal $0$-dof-graph, as in Case~III.

\begin{lemma}[The graph $G - v$ is a minimal $k'$-dof-graph with $k' \le D - 2$; from the proof of KT Claim~6.11]
  \label{lem:case-III-gap3-minimalKDof}
  \lean{Graph.splitOff_removeVertex_minimalKDof}
  \leanok
  \uses{def:k-dof, def:graph-operations, lem:subgraph-minimality, thm:def-eq-corank,
        lem:splitoff-deficiency, lem:case-III-claim-6-11-base}
  Let $D \ge 2$, and let $G$ be a minimal $0$-dof-graph with a vertex $v$
  of degree two, whose edges $e_a$ and $e_b$ join $v$ to distinct
  vertices $a, b \neq v$. Let $G_v = G - v$ and
  $k' = \operatorname{def}(\tilde G_v)$. Then $G_v$ is a minimal
  $k'$-dof-graph, and
  \[
    0 \;\le\; k' \;\le\; D - 2.
  \]
\end{lemma}
\begin{proof}
  \leanok
  We follow Katoh and Tanigawa~\cite[p.~684]{katohTanigawa2011}. The
  graph $G_v$ is a subgraph of $G$, so it is a minimal $k'$-dof-graph by
  KT Lemma~3.3 (\cref{lem:subgraph-minimality}). It has the vertex $a$,
  so the partition of $V(G_v)$ into one part shows that $k' \ge 0$
  (\cref{def:D-deficiency}). For the same reason
  $\operatorname{def}(\widetilde{G_v^{ab}}) \ge 0$. By KT Lemma~4.3(i)
  (\cref{lem:splitoff-deficiency}),
  $\operatorname{def}(\widetilde{G_v^{ab}}) \le \operatorname{def}(\tilde G) = 0$,
  so $G_v^{ab}$ is a $0$-dof-graph.

  By \cref{lem:case-III-claim-6-11-base}, at $k = 0$,
  $M(\widetilde{G_v^{ab}})$ has a base $B'$ with
  $h = |\widetilde{ab} \cap B'| < D - 1$. As $G_v^{ab}$ has the
  $|V(G)| - 1$ vertices of $G$ other than $v$, KT eq.~(2.4)
  (\cref{thm:def-eq-corank}) gives $|B'| = D(|V(G)| - 2)$. The graph
  $G_v$ is $G_v^{ab}$ without the edge $ab$
  (\cref{lem:splitoff-edge-substitution}). So $M(\tilde G_v)$ is the
  restriction of $M(\widetilde{G_v^{ab}})$ to $E(\tilde G_v)$
  (\cref{lem:matroid-restrict-subgraph}), in which
  $B' \setminus \widetilde{ab}$ is an independent set of $|B'| - h$
  elements. As $G_v$ and $G_v^{ab}$ have the same vertices, KT
  eq.~(2.4) gives
  \[
    k' \;=\; D\,(|V(G)| - 2) - \operatorname{rank} M(\tilde G_v)
    \;\le\; |B'| - \bigl(|B'| - h\bigr) \;=\; h \;\le\; D - 2.
  \]
\end{proof}

\begin{fmlnote}
  The pinned theorem also assumes that some edge label lies outside
  $E(G)$, the label of the new edge $ab$ of $G_v^{ab}$, which the proof
  uses (\cref{rem:fresh-edge-supply}).
\end{fmlnote}

Katoh and Tanigawa prove KT eq.~(6.22) at the normals $q$ of KT
eq.~(6.18), from the genericity of $q$ (the remark below). As
explained after \cref{lem:case-I-realization}, we replace genericity by
the non-vanishing of polynomials in the normals. So we state KT
eq.~(6.22) as a bound off the zeros of a polynomial.

\begin{lemma}[The rank of $G - v$ off the zeros of a polynomial; KT eq.~(6.22)]
  \label{lem:case-III-nested-rank-lower}
  \lean{CombinatorialRigidity.Molecular.PanelHingeFramework.exists_nested_rankPolynomial_lower_all_k}
  \leanok
  \uses{def:rank-hypothesis, lem:case-III-gap3-minimalKDof}
  Let $d \ge 2$, let $G$ be a minimal $0$-dof-graph with a vertex $v$ of
  degree two, whose edges join $v$ to distinct vertices $a, b \neq v$,
  and let $G_v = G - v$. Then there is a nonzero polynomial $P$ in the
  normals such that
  \[
    \operatorname{rank} R(G_v, q) \;\ge\; D\,(|V(G_v)| - 1) - (D - 2)
  \]
  for all normals $q$ with $P(q) \neq 0$.
\end{lemma}
\begin{proof}
  \leanok
  We assume, as in Case~III, that $G$ is simple and that the induction
  hypothesis KT eq.~(6.1) holds for the graphs with fewer vertices than
  $G$ (the note after the proof). By
  \cref{lem:case-III-gap3-minimalKDof}, $G_v$ is a minimal
  $k'$-dof-graph with $k' = \operatorname{def}(\tilde G_v) \le D - 2$.
  It is simple, as a subgraph of $G$, and it has at least one and fewer
  than $|V(G)|$ vertices. So the induction hypothesis gives a generic
  full-rank realization $(G_v, q_0)$ (\cref{def:rank-hypothesis}), with
  \[
    \operatorname{rank} R(G_v, q_0) \;=\; D\,(|V(G_v)| - 1) - k'.
  \]
  Its normals are in general position and $G_v$ has no loops, so its
  hinges are nondegenerate (\cref{lem:general-position-support-nonzero}).
  By \cref{lem:rank-polynomial-of-le-finrank}, there is a polynomial $P$
  in the normals with $P(q_0) \neq 0$ such that
  $\operatorname{rank} R(G_v, q) \ge D(|V(G_v)| - 1) - k'$ wherever
  $P(q) \neq 0$. In particular $P \neq 0$. As $k' \le D - 2$, this is
  the bound of the lemma.
\end{proof}

\begin{fmlnote}
  The pinned theorem assumes that $G$ is simple, and the induction
  hypothesis for every $k'$ and every minimal $k'$-dof-graph with at
  least one and fewer than $|V(G)|$ vertices. The statement omits both.
  The pinned theorem also assumes that some edge label lies outside
  $E(G)$ (\cref{rem:fresh-edge-supply}). It fixes an endpoint selector
  that assigns each edge of $G_v^{ab}$ its ends
  (\cref{fmlnote:panel-hinge-nondegenerate}) before it gives $P$.
\end{fmlnote}

The lemma needs $G$ to be simple. Let $d = 2$, so that $D = 3$, and let
$G$ have the vertices $v, a, b, c$, the edges $va$, $vb$ and $cb$, and
two edges $ac$. Comparing the partitions of $V(G)$
(\cref{def:D-deficiency}) shows that $\operatorname{def}(\tilde G) = 0$.
Removing an edge $ac$ gives the value $1$ to the partition into
singletons. Removing $va$ or $vb$ gives it to the partition
$\{\{v\}, \{a, b, c\}\}$, and removing $cb$ to
$\{\{b\}, \{a, c, v\}\}$. So $G$ is a minimal $0$-dof-graph, and $v$ has
degree two. In $(G_v, q)$ the two edges $ac$ have the same hinge, and so
the same rows. Let $Q$ be the product of a $2 \times 2$ minor of the
normals at $a$ and $c$ and one of the normals at $c$ and $b$. Where
$Q(q) \neq 0$, the hinges of $(G_v, q)$ are nondegenerate, and
\[
  \operatorname{rank} R(G_v, q) \;\le\; 2(D - 1) \;=\; 4
  \;<\; 5 \;=\; D\,(|V(G_v)| - 1) - (D - 2).
\]
Over an infinite field, $PQ$ has a non-root for every nonzero
polynomial $P$, so no polynomial $P$ as in the lemma exists. The two
edges $ac$ form a proper rigid subgraph of $G$, so this graph falls under
Case~I.

\begin{remark}[KT eq.~(6.22) and KT footnote~6]
  We recall Katoh and Tanigawa's proof of KT
  eq.~(6.22)~\cite[pp.~684--685]{katohTanigawa2011}. The realization
  $(G_v^{ab}, q)$ of KT eq.~(6.18) is generic: the coordinates of the
  vectors $c(u)$ that define its panels are algebraically independent
  over $\Q$ (KT \S5.1, p.~668; the passage after
  \cref{lem:case-I-realization}). By KT
  eq.~(6.1), the simple minimal $k'$-dof-graph $G_v$ has a nonparallel
  realization of rank $D(|V(G_v)| - 1) - k'$. A generic nonparallel
  realization of $G_v$ has the largest rank of all nonparallel
  realizations of $G_v$ (KT \S5.1). No realization with nondegenerate
  hinges has a rank larger than $D(|V(G_v)| - 1) - k'$
  (\cref{lem:relative-deficiency-rank-bound}). So every
  generic nonparallel realization of $G_v$ has rank
  $D(|V(G_v)| - 1) - k'$ (KT footnote~6, p.~684). The realization
  $(G_v, q)$ has the bodies and the panels of $(G_v^{ab}, q)$, so it is a
  generic nonparallel realization of $G_v$. Hence
  $\operatorname{rank} R(G_v, q) = D(|V(G_v)| - 1) - k'$.

  In our proof the normals of KT eq.~(6.18) are not given first. Case~III
  chooses them off the zeros of a product of nonzero polynomials, fixed
  before the choice (the proof of \cref{lem:case-III-chain-discriminator},
  later in this subsection). One factor is the polynomial $P$ of
  \cref{lem:case-III-nested-rank-lower}. Another is the polynomial that
  \cref{lem:rank-polynomial-of-le-finrank} gives for the realization of
  $G_v^{ab}$ from the induction hypothesis, and a third puts the normals
  in general position. Off the zeros of the second,
  $\operatorname{rank} R(G_v^{ab}, q) \ge D(|V(G_v)| - 1)$. Off those of
  the third, the hinges are nondegenerate, and as
  $\operatorname{def}(\widetilde{G_v^{ab}}) = 0$, the rank is at most
  $D(|V(G_v)| - 1)$ (\cref{lem:relative-deficiency-rank-bound}). So KT
  eqs.~(6.18) and~(6.22) hold at the same normals, as in Katoh and
  Tanigawa's proof. They hold by the choice of the normals, not by their
  genericity.
\end{remark}

The third step is a count in linear algebra, which we state for one
hinge of any body-hinge framework.

\begin{lemma}[A redundant row at a hinge; from the proof of KT Claim~6.11]
  \label{lem:case-III-claim-6-11-redundant-row}
  \lean{CombinatorialRigidity.Molecular.BodyHingeFramework.exists_redundant_panelRow_of_edge_of_finrank_lt}
  \leanok
  \uses{lem:case-III-gap3-minimalKDof, lem:case-II-placement-new-rows}
  Let $(G, p)$ be a body-hinge framework, let $e$ be an edge of $G$ with
  distinct ends whose hinge is nondegenerate, and let $R_e$ be the span
  of the rows of $R(G, p)$ at $e$. Let $W$ be a space of linear forms on
  the screw assignments such that
  \[
    \dim (W + R_e) \;<\; \dim W + (D - 1).
  \]
  Then $R_e$ has a basis $r_1, \dots, r_{D-1}$ with
  $r_i \in W + \operatorname{span}\{r_j : j \neq i\}$ for some $i$.
\end{lemma}
\begin{proof}
  \leanok
  By \cref{lem:case-II-placement-new-rows}, $R_e$ contains $D - 1$
  linearly independent forms $r_1, \dots, r_{D-1}$. The space $R_e$ has
  dimension $D - 1$, as in the proof of that lemma, so they form a basis
  of $R_e$. Suppose that no $r_i$ lies in
  $W + \operatorname{span}\{r_j : j \neq i\}$. If
  $\sum_j c_j r_j \in W$ with $c_i \neq 0$, then dividing by $c_i$
  puts $r_i$ in $W + \operatorname{span}\{r_j : j \neq i\}$. So
  $\sum_j c_j r_j \in W$ only if every $c_j$ is $0$, and the images of
  the $r_j$ in the quotient by $W$ are linearly independent. Then
  $\dim (W + R_e) = \dim W + (D - 1)$, contrary to the hypothesis.
\end{proof}

\begin{fmlnote}
  As for \cref{lem:case-II-placement-new-rows}, the pinned theorem holds
  for every body-hinge framework, and it reads the ends of $e$ from the
  endpoint selector (\cref{fmlnote:panel-hinge-nondegenerate}).
\end{fmlnote}

We apply this count at the edge $ab$ of $G_v^{ab}$, where $W$ is the row
space of $R(G_v, q)$.

\begin{lemma}[KT Claim~6.11: a redundant row at $ab$; KT eq.~(6.23)]
  \label{lem:case-III-claim-6-11}
  \lean{CombinatorialRigidity.Molecular.BodyHingeFramework.exists_redundant_panelRow_ab_of_finrank_eq}
  \leanok
  \uses{lem:case-III-claim-6-11-redundant-row, lem:case-III-gap3-minimalKDof}
  Let $d \ge 2$, let $G$ be a multigraph with a vertex $v$ of degree
  two, whose edges join $v$ to distinct vertices $a, b \neq v$, and let
  $G_v = G - v$. Let $q$ be normals at which the hinge of the new edge
  $ab$ of $(G_v^{ab}, q)$ is nondegenerate. Suppose that, for some
  $k' \le D - 2$,
  \[
    \operatorname{rank} R(G_v^{ab}, q) = D\,(|V(G_v)| - 1)
    \quad\text{and}\quad
    \operatorname{rank} R(G_v, q) = D\,(|V(G_v)| - 1) - k',
  \]
  as in KT eqs.~(6.18) and~(6.22). Then one of the $D - 1$ rows of
  $R(G_v^{ab}, q)$ at $ab$ lies in the span of the rows of $R(G_v, q)$
  and the other rows at $ab$.
\end{lemma}
\begin{proof}
  \leanok
  We follow Katoh and Tanigawa~\cite[p.~685]{katohTanigawa2011}. The
  edges of $G_v^{ab}$ are those of $G_v$, with the same ends, and the
  new edge $ab$ (\cref{lem:splitoff-edge-substitution}). The two
  frameworks have the same normals, so an edge of $G_v$ has the same
  hinge, and the same rows, in both. Let $W$ be the row space of
  $R(G_v, q)$, and let $R_{ab}$ be the span of the rows of
  $R(G_v^{ab}, q)$ at $ab$. Then the row space of $R(G_v^{ab}, q)$ is
  $W + R_{ab}$, and the two rank equations give
  \[
    \dim (W + R_{ab}) \;=\; D\,(|V(G_v)| - 1) \;=\; \dim W + k'.
  \]
  As $k' \le D - 2 < D - 1$, the rows at $ab$ raise $\dim W$ by fewer
  than $D - 1$.

  The $D - 1$ rows $r_1, \dots, r_{D-1}$ of $R(G_v^{ab}, q)$ at $ab$
  form a basis of $R_{ab}$ (\cref{def:hinge-row-block}). The proof of
  \cref{lem:case-III-claim-6-11-redundant-row} applies to every basis of
  $R_{ab}$. So some $r_i$ lies in
  $W + \operatorname{span}\{r_j : j \neq i\}$.
  The rows of $R(G_v, q)$ and the $r_j$ with $j \neq i$ then span
  $W + R_{ab}$. So removing the row $r_i$ does not lower the rank of
  $R(G_v^{ab}, q)$, which is KT eq.~(6.23).
\end{proof}

\begin{fmlnote}
  The pinned theorem gives a basis of the span of the rows at $ab$, one
  member of which lies in the span of the rows of $R(G_v, q)$ and the
  other members. This is the lemma for one choice of the rows at $ab$
  (\cref{def:hinge-row-block}). The pinned theorem holds for any two
  multigraphs whose edges, with their ends, agree except for one edge
  $ab$ of the first, and with any number in place of $|V(G_v)|$. Their
  frameworks share one endpoint selector, which assigns $ab$ its ends
  (\cref{fmlnote:panel-hinge-nondegenerate}). The formal proof of
  Case~III applies it with
  $k' = D(|V(G_v)| - 1) - \operatorname{rank} R(G_v, q)$, for which the
  second rank equation holds by definition, and
  \cref{lem:case-III-nested-rank-lower} gives $k' \le D - 2$.
\end{fmlnote}

\subsubsection{The candidate-completion: the redundant row becomes the missing row
(KT eqs.\ (6.24)--(6.29))}
\label{sec:molecular-algebraic-induction-candidate-completion}

Claim~6.11 (\cref{lem:case-III-claim-6-11}) produces a redundant $ab$-row of
$R(G_v^{ab}, q)$; the candidate-completion converts it into the one row that lifts the
degenerate placement of \cref{lem:case-II-realization-placement} from rank $D(|V| - 1) - 1$
to full rank $D(|V| - 1)$. KT's mechanism (eqs.~(6.24)--(6.29)) is a single row operation
followed by the eqs.~(6.14)--(6.15) column operation. The redundant row is first written as
an explicit vanishing combination of rigidity rows (eqs.~(6.24)--(6.25),
\cref{lem:case-III-redundant-decomposition}); transported from $R(G_v^{ab}, q)$ to the
degenerate placement $R(G, p_1)$ across the seam identity (\cref{lem:case-III-seam}) and
collapsed by the eq.~(6.27) hinge-difference identity
(\cref{lem:case-III-transport-collapse}, \cref{lem:case-III-candidate-row-construction}), it
becomes a row $w$ of $R(G, p_1)$ supported on both bodies $v$ and $a$. The column operation
$\mathrm{col}_a \mathrel{+}= \mathrm{col}_v$ (\cref{lem:case-III-columnop}) then leaves $w$
depending only on body $v$'s screw column (eq.~(6.28), \cref{lem:case-III-vanish-off-column}),
so it lifts the $va$-block from $D - 1$ rows to $D$ (\cref{lem:case-III-conditional-block}).
Transporting independence back through the column operation gives a full-rank family of size
$((D - 1) + 1) + D(|V_v| - 1) = D(|V| - 1)$ (\cref{lem:case-III-candidate-row}). The assembly
passes through one conditional --- the eq.~(6.29) top-left $D \times D$ block has full rank for
some choice of line --- discharged below by Claim~6.12 (\cref{lem:case-III-claim612}).

\begin{lemma}[A rigidity row vanishing off body $v$'s column is a pure $v$-column row; KT eq.\ (6.28)]
  \label{lem:case-III-vanish-off-column}
  \lean{CombinatorialRigidity.Molecular.BodyHingeFramework.dualMap_eq_comp_single_proj_of_vanish_off}
  \leanok
  Let $f$ be a linear functional on the screw assignments $S \colon \alpha \to
  \bigwedge^k K^{k+2}$ that vanishes whenever $S_v = 0$. Then $f$ depends only on body $v$'s
  screw column: $f(S) = f_v(S_v)$, where $f_v$ is the restriction of $f$ to that column.
\end{lemma}
\begin{proof}
  \leanok
  Write $S = \iota_v(S_v) + (S - \iota_v(S_v))$, where $\iota_v$ places a screw on body $v$ and
  zero elsewhere. The second summand vanishes at $v$, so $f$ kills it by hypothesis, leaving
  $f(S) = f(\iota_v(S_v)) = f_v(S_v)$.
\end{proof}

\begin{lemma}[The column operation $\Phi = \mathrm{col}_a \mathrel{+}= \mathrm{col}_v$; KT eqs.\ (6.14)--(6.16)]
  \label{lem:case-III-columnop}
  \lean{CombinatorialRigidity.Molecular.BodyHingeFramework.columnOp,
        CombinatorialRigidity.Molecular.BodyHingeFramework.hingeRow_comp_columnOp_apply,
        CombinatorialRigidity.Molecular.BodyHingeFramework.hingeRow_comp_columnOp_vanish_off}
  \leanok
  \uses{lem:case-III-vanish-off-column}
  Fix distinct bodies $v \ne a$. KT's column operation $\mathrm{col}_a \mathrel{+}= \mathrm{col}_v$
  (eqs.~(6.14)--(6.15)), acting on the rigidity rows through the relative screws $S_u - S_w$ they
  read, is the change of variables $\Phi$ on screw assignments that replaces $S_v$ by $S_v + S_a$ and
  fixes every other body. It is a linear automorphism (inverse: replace $S_v$ by $S_v - S_a$), hence
  rank-preserving. For a hinge row $S \mapsto \rho(S_v - S_a)$ --- supported on both bodies $v$ and
  $a$ --- the operated row is pure $v$-column:
  \[ \rho(S_v - S_a) \circ \Phi = \rho\bigl((S_v + S_a) - S_a\bigr) = \rho(S_v), \]
  which vanishes whenever $S_v = 0$.
\end{lemma}
\begin{proof}
  \leanok
  $\Phi$ is a linear automorphism by direct verification of additivity, homogeneity, and the two
  inverse identities (a case split on whether a coordinate is $v$: the off-$v$ coordinates are
  untouched, and at $v$ the two composites cancel by $(S_v + S_a) - S_a = S_v$ and $(S_v - S_a) + S_a
  = S_v$, using $v \ne a$ so $\Phi(S)_a = S_a$). The evaluation $\rho((S_v + S_a) - S_a) = \rho(S_v)$
  is then immediate from $\Phi(S)_v = S_v + S_a$ and $\Phi(S)_a = S_a$; when $S_v = 0$ it reads
  $\rho(0) = 0$.
\end{proof}

\begin{lemma}[The seam identity: an edge's row depends only on the seed at its endpoints; KT eq.\ (6.26)]
  \label{lem:case-III-seam}
  \lean{CombinatorialRigidity.Molecular.PanelHingeFramework.ofNormals_panelRow_eq_of_ends_seed_eq}
  \leanok
  Let two free-normal panel realizations $R(G, q_0)$ and $R(G_v^{ab}, q)$, on possibly different
  graphs but sharing one endpoint selector, and let $e$ be an edge whose two endpoints carry the same
  normals under $q_0$ and $q$. Then $e$ contributes the same rigidity rows to both: for every index
  pair $(t_1, t_2)$,
  \[ R(G, q_0; e, t_1, t_2) = R(G_v^{ab}, q; e, t_1, t_2). \]
\end{lemma}
\begin{proof}
  \leanok
  A panel row at $e$ is determined by the two endpoints of $e$, the row index, and the supporting
  extensor of the two endpoint normals; the underlying graph never enters. Agreement of $q_0$ and $q$
  at the two endpoints makes the endpoint normals equal, hence the supporting extensors and the rows.
\end{proof}

\begin{lemma}[The redundant $ab$-row as an explicit vanishing combination; KT eqs.\ (6.24)--(6.25)]
  \label{lem:case-III-redundant-decomposition}
  \lean{CombinatorialRigidity.Molecular.BodyHingeFramework.exists_redundant_panelRow_ab_decomposition}
  \leanok
  \uses{lem:case-III-claim-6-11}
  Let $r$ be the $D - 1$ independent $ab$-rows of \cref{lem:case-III-claim-6-11}, with $r_{i^*}$ the
  redundant one. Then $r_{i^*}$ decomposes as
  \[
    r_{i^*} = w_{\mathrm{Gv}} + w_{\mathrm{Other}}, \qquad
    w_{\mathrm{Gv}} \in \operatorname{span}\bigl(R(G_v, q)\text{-rows}\bigr), \quad
    w_{\mathrm{Other}} \in \operatorname{span}\{r_j : j \ne i^*\},
  \]
  KT's eq.~(6.24): the redundant row is a combination of the $G_v$-rows and the other $ab$-rows.
  Equivalently, the vanishing combination $r_{i^*} - w_{\mathrm{Other}} - w_{\mathrm{Gv}} = 0$ has
  $(ab)i^*$-coefficient $\lambda_{(ab)i^*} = 1$ (eq.~(6.25)).
\end{lemma}
\begin{proof}
  \leanok
  Claim~6.11 (\cref{lem:case-III-claim-6-11}) places $r_{i^*}$ in the sum of the two spans; split
  that membership into its $\operatorname{span}(R(G_v, q)\text{-rows})$ summand $w_{\mathrm{Gv}}$ and
  its $\operatorname{span}\{r_j : j \ne i^*\}$ summand $w_{\mathrm{Other}}$.
\end{proof}

The following restriction of the eq.~(6.24) combination is not needed for the candidate row itself;
it is the input to the $M_3$ candidate of Claim~6.12, below.

\begin{lemma}[The $a$-column block of the eq.\ (6.24) vanishing combination is $0$; KT eq.\ (6.43)]
  \label{lem:case-III-acolumn-zero}
  \lean{CombinatorialRigidity.Molecular.BodyHingeFramework.exists_redundant_panelRow_ab_decomposition_acolumn_zero}
  \leanok
  \uses{lem:case-III-redundant-decomposition}
  The eqs.~(6.24)--(6.25) decomposition (\cref{lem:case-III-redundant-decomposition}) makes
  $g := w_{\mathrm{Gv}} + w_{\mathrm{Other}} - r_{i^*}$ the zero functional --- KT's eq.~(6.24)
  vanishing combination
  \[
    \sum_j \lambda_{(ab)j}\, R(G_v^{ab}, q; (ab)j) + \sum_{e \in E_v, j} \lambda_{ej}\,
      R(G_v^{ab}, q; e_j) = 0
  \]
  on screw assignments. Its restriction to any single body $a$'s screw column vanishes as well
  (KT eq.~(6.43)):
  \[ \sum_{e \in E_v \cup \{ab\},\, j} \lambda_{ej}\, R(G_v^{ab}, q; e_j, a) = 0. \]
\end{lemma}
\begin{proof}
  \leanok
  By \cref{lem:case-III-redundant-decomposition} $r_{i^*} = w_{\mathrm{Gv}} + w_{\mathrm{Other}}$, so
  $g$ is the zero functional; restricting the zero functional to body $a$'s screw column is $0$ for
  every $a$.
\end{proof}

\begin{lemma}[The conditional $D$-row block: the operated row lifts the $va$-block from $D-1$ to $D$; KT eq.\ (6.29)]
  \label{lem:case-III-conditional-block}
  \lean{CombinatorialRigidity.Molecular.BodyHingeFramework.linearIndependent_sum_pinned_block_augment}
  \leanok
  \uses{lem:case-II-placement-block-independent, lem:case-III-vanish-off-column,
        lem:case-III-columnop}
  KT eq.~(6.29)'s block-triangular assembly, taking the rank $D(|V|-1) - 1$ of
  \cref{lem:case-II-realization-placement} to full $D(|V|-1)$, conditional on the top-left
  $D \times D$ block being full rank. Fix a body $v$. Take a new block $r_n$ of $D - 1$ rows
  together with one further row $w$, and an old block $r_o$, linearly independent, each of whose
  rows vanishes on every screw assignment supported on body $v$ alone. If the restrictions of
  $r_n$ and $w$ to body $v$'s $D$-dimensional screw column form a linearly independent family ---
  KT's eq.~(6.29) top-left $D \times D$ full-rank condition --- then the combined family
  $r_n \cup \{w\} \cup r_o$ is linearly independent. The count is
  $((D-1)+1) + D(|V_v|-1) = D(|V|-1)$.

  In the candidate-completion, $r_n$ is body $v$'s $va$-block and $w$ is the operated candidate row
  (\cref{lem:case-III-columnop}, \cref{lem:case-III-vanish-off-column}); $r_o$ is the split-off
  $G_v^{ab} \setminus (ab)i^*$.
\end{lemma}
\begin{proof}
  \leanok
  \uses{lem:case-II-placement-block-independent}
  This is the pin-a-body column split (\cref{lem:case-II-placement-block-independent}) applied to
  the new block augmented by $w$: restricting the augmented block to body $v$'s column reproduces
  the hypothesized independent family, and the old block's vanishing on assignments supported at
  $v$ lets it join without further constraint.
\end{proof}

\begin{lemma}[The eq.\ (6.27) transport collapse: the $(vb)$-row minus the $(ab)$-row is a $(va)$-hinge row; KT eqs.\ (6.26)--(6.27)]
  \label{lem:case-III-transport-collapse}
  \lean{CombinatorialRigidity.Molecular.PanelHingeFramework.panelRow_vb_sub_panelRow_ab_eq_hingeRow_va,
        CombinatorialRigidity.Molecular.BodyHingeFramework.hingeRow_sub_hingeRow_eq,
        CombinatorialRigidity.Molecular.BodyHingeFramework.comp_columnOp_eq_comp_single_proj}
  \leanok
  \uses{lem:case-III-seam, lem:case-III-columnop, lem:case-III-vanish-off-column}
  At the degenerate eq.~(6.12) placement $p_1$ --- body $v$'s normal set to $n_a + t\,n_b$, agreeing
  with $q$ at the surviving endpoint $b$ --- the transported $(vb)j$-row of $R(G, p_1)$ and the
  inductive $(ab)j$-row of $R(G_v^{ab}, q)$ read the same supporting extensor $C = C(n_a, n_b)$: the
  shear identity $C(n_a + t\,n_b, n_b) = C(n_a, n_b)$ (KT eq.~(6.16)). Writing the two rows as
  $\rho(S_v - S_b)$ and $\rho(S_a - S_b)$ for the common functional $\rho$, their difference is the
  pure $(va)$-hinge row
  \[ R(G, p_1; (vb)j) - R(G_v^{ab}, q; (ab)j) = \rho(S_v - S_a), \]
  since $(S_v - S_b) - (S_a - S_b) = S_v - S_a$. This is KT eq.~(6.27): transporting the redundant
  combination as $(vb)$-rows and subtracting the inductive $(ab)$-combination --- which the
  eq.~(6.24) decomposition makes vanish --- collapses the transported row to a single $(va)$-hinge
  row $w$. A companion identity records that the operated row $w \circ \Phi$
  (\cref{lem:case-III-columnop}) equals its own pure $v$-column factorization
  (\cref{lem:case-III-vanish-off-column}), the eq.~(6.29) $(vb)i^*$-row.
\end{lemma}
\begin{proof}
  \leanok
  Both rows read the panel through $b$ with the same supporting extensor $C$ --- the shear identity
  at $p_1$ against the $(ab)$-extensor at $q$ --- so they are $\rho(S_v - S_b)$ and $\rho(S_a - S_b)$;
  their difference is $\rho(S_v - S_a)$. The companion identity composes
  \cref{lem:case-III-columnop} (the operated row vanishes off $v$) with
  \cref{lem:case-III-vanish-off-column} (such a row factors through $v$'s column).
\end{proof}

\begin{lemma}[The candidate row $w$ from the redundant row's common element; KT eqs.\ (6.24)--(6.28)]
  \label{lem:case-III-candidate-row-construction}
  \lean{CombinatorialRigidity.Molecular.PanelHingeFramework.exists_candidate_row_eq612}
  \leanok
  \uses{lem:case-III-redundant-decomposition, lem:case-III-transport-collapse}
  The common element $w_{\mathrm{Gv}}$ of the eqs.~(6.24)--(6.25) decomposition
  (\cref{lem:case-III-redundant-decomposition}) --- the $G_v$-row part of the redundant $ab$-row ---
  is itself an $ab$-block element (it equals $r_{i^*} - w_{\mathrm{Other}}$, both terms in the block).
  The $ab$-block is the image, under the $ab$-hinge, of the $(D-1)$-dimensional hinge-row block
  $r(p(e_0)) = (\operatorname{span} C)^\perp$, so $w_{\mathrm{Gv}} = \rho(S_a - S_b)$ for some
  $\rho \in r(p(e_0))$. At the eq.~(6.12) seed the shear identity $C(n_a + t\,n_b, n_b) = C(n_a, n_b)$
  (KT eq.~(6.16)) makes $\rho$ a hinge-row-block functional of the $vb$-hinge of $R(G, p_1)$ too,
  so $\rho(S_v - S_b)$ is the transported $(vb)i^*$-row, a rigidity row of $R(G, p_1)$. Its
  eq.~(6.27) collapse against the inductive $(ab)$-part is the pure $(va)$-hinge candidate row
  \[ \rho(S_v - S_b) \;-\; w_{\mathrm{Gv}} \;=\; \rho(S_v - S_a) \;=\; w \]
  (\cref{lem:case-III-transport-collapse}).
\end{lemma}
\begin{proof}
  \leanok
  \uses{lem:case-III-transport-collapse}
  The $ab$-block is the image of the hinge-row block $r(p(e_0))$ under the $ab$-hinge, so
  $w_{\mathrm{Gv}} = \rho(S_a - S_b)$ for some $\rho \in r(p(e_0))$. The shear identity gives
  $C(p_1(e_b)) = C(q(e_0))$, so the same $\rho$ annihilates the $vb$-hinge's supporting extensor and
  $\rho(S_v - S_b)$ is a rigidity row of $R(G, p_1)$ at the $vb$-link. Subtracting $w_{\mathrm{Gv}}$
  collapses it to $\rho(S_v - S_a) = w$ by the hinge-difference identity
  (\cref{lem:case-III-transport-collapse}).
\end{proof}

\begin{lemma}[The candidate-completion row operation: the missing row $w$; KT eqs.\ (6.24)--(6.29)]
  \label{lem:case-III-candidate-row}
  \lean{CombinatorialRigidity.Molecular.BodyHingeFramework.linearIndependent_sum_augment_candidateRow}
  \leanok
  \uses{lem:case-III-claim-6-11, lem:case-III-redundant-decomposition,
        lem:case-III-candidate-row-construction,
        lem:case-III-vanish-off-column, lem:case-III-columnop, lem:case-III-seam,
        lem:case-III-transport-collapse,
        lem:case-III-conditional-block, lem:case-III-claim612,
        lem:case-II-realization-placement}
  Realize $G$ at the degenerate eq.~(6.12) placement $p_1$ with $p_1(vb) = q(ab)$
  (\cref{lem:case-II-realization-placement}). Provided the eq.~(6.29) top-left $D \times D$ block has
  full rank for some choice of line (Claim~6.12, \cref{lem:case-III-claim612}), the candidate row $w$
  of \cref{lem:case-III-candidate-row-construction}, operated by the column operation
  (\cref{lem:case-III-columnop}), joins the $va$-block and lifts $R(G, p_1)$ to a full-rank family of
  size $D(|V| - 1)$: the matrix becomes block-triangular as in KT eq.~(6.29), with bottom-right
  $R(G_v^{ab} \setminus (ab)i^*, q)$ of rank $D(|V \setminus \{v\}| - 1)$.
\end{lemma}
\begin{proof}
  \leanok
  \uses{lem:case-III-candidate-row-construction, lem:case-III-columnop,
        lem:case-III-vanish-off-column, lem:case-III-transport-collapse,
        lem:case-III-conditional-block, lem:case-III-claim612,
        lem:case-II-realization-placement}
  The redundant $ab$-row, decomposed by \cref{lem:case-III-redundant-decomposition} into a $G_v$-row
  part and an other-$ab$-row part with $\lambda_{(ab)i^*} = 1$, transports up to $R(G, p_1)$ by the
  eq.~(6.26) row correspondence (the $(ab)j$-rows of $R(G_v^{ab}, q)$ become the $(vb)j$-rows of
  $R(G, p_1)$; the $E_v$-rows correspond directly). \Cref{lem:case-III-candidate-row-construction}
  reads the common $G_v$-row element as $\rho(S_a - S_b)$ for a hinge-row-block functional $\rho$,
  identifies its transported $(vb)i^*$-row $\rho(S_v - S_b)$ as a rigidity row of $R(G, p_1)$, and
  collapses it to the candidate row $w = \rho(S_v - S_a)$ (\cref{lem:case-III-transport-collapse}).

  The candidate row is supported on both bodies $v$ and $a$. The column operation
  $\mathrm{col}_a \mathrel{+}= \mathrm{col}_v$ (\cref{lem:case-III-columnop}) leaves $w \circ \Phi$
  depending only on body $v$'s column (\cref{lem:case-III-vanish-off-column}), so its
  $V \setminus \{v\}$ part vanishes (eq.~(6.28)). Precomposing the whole family $r_n \cup \{w\} \cup
  r_o$ with the automorphism $\Phi$ preserves independence; in the operated frame $w \circ \Phi$
  joins the $va$-block in the pin-a-body augment (\cref{lem:case-III-conditional-block}), while the
  old block is fixed (since $\Phi$ fixes bodies other than $v$) and keeps its vanishing. The augment
  is independent provided the eq.~(6.29) top-left $D \times D$ block --- the $D - 1$ pinned $va$-rows
  plus the operated $w$'s $v$-column --- has full rank; this is the one conditional, discharged by
  Claim~6.12 (\cref{lem:case-III-claim612}). Transporting the operated independence back through
  $\Phi$ gives the original family independent, of size $((D-1)+1) + D(|V_v|-1) = D(|V|-1)$, with
  bottom-right $R(G_v^{ab} \setminus (ab)i^*, q)$ of rank $D(|V_v|-1)$. The eq.~(6.43) $a$-column
  vanishing (\cref{lem:case-III-acolumn-zero}) is a Claim~6.12 fact (the $M_3$ case), not used here.
\end{proof}

\subsubsection{The three candidates and Claim 6.12: the \texorpdfstring{$d = 3$}{d = 3} worked case (KT \S6.4.1, eqs.\ (6.30)--(6.45))}
\label{sec:molecular-algebraic-induction-claim612}

The remainder of this subsection works out Katoh--Tanigawa's three-dimensional case
(KT \S6.4.1, Lemma~6.10) in full, as the accessible instance of the general argument. The
general Lemma~6.13 (\cref{lem:case-III}, KT \S6.4.2) applies the same three-candidate idea
along a length-$d$ chain, but at $d = 3$ each of its moving parts reduces to something
concrete, so the three-dimensional case carries the whole idea --- three candidate
placements, one of them full rank by a six-join span --- without the chain, dichotomy, and
block scaffolding the general form otherwise asks the reader to absorb at once.

Several reductions make the $d = 3$ case self-contained. The redundant $ab$-row is tested
against three \emph{fixed} candidates by a case split over three panels, rather than by a
discriminator ranging over the interior vertices of a length-$d$ chain. The third candidate
is reached by a single vertex relabel exchanging $v$ and $a$
(\cref{lem:splitOff-ofNormals-relabel}), rather than by transporting the redundancy along
the chain one column operation per interior vertex. Each candidate is certified by a direct
three-panel rank count, rather than by the block certificate with the moving-vertex
bookkeeping the general chain needs. The base triangle supplies the three candidate panels
directly, so the chain-versus-short-cycle dichotomy (KT Lemmas~4.6/4.8) and the cycle family
(KT Lemma~5.4) --- present in the general proof --- play no part at $d = 3$. And the closing
span argument runs on the six joins of four points in $\bigwedge^2 K^4$, a six-dimensional
space small enough to picture, rather than on the $\binom{k+2}{2}$ joins of
$\bigwedge^k K^{k+2}$.

The candidate-completion above lifts a
\emph{single} degenerate placement $p_1$ to full rank $D(|V|-1)$ \emph{provided} its top-left
$D \times D$ block (eq.~(6.29)) is full rank for some line choice (Claim~6.12,
\cref{lem:case-III-claim612}). At $d = 2$ KT
show this holds for $p_1$ or its symmetric partner $p_2$; \emph{at $d = 3$ it may fail for both}, so
the conditional is genuinely an existential over a free choice of lines. KT introduce a third candidate $p_3$ ---
splitting off at the \emph{other} degree-2 vertex $a$ along $vc$, whose split-off graph $G_a^{vc}$ is
isomorphic to $G_v^{ab}$ (eqs.~(6.31)--(6.41)) --- and prove Claim~6.12: at least one of the three
top-left blocks $M_1, M_2, M_3$ (of eqs.~(6.29), (6.30), (6.41)) has full rank, for some choice of
lines $L \subset \Pi(a)$, $L' \subset \Pi(b)$, $L'' \subset \Pi(c)$ (writing $\Pi(u) := \Pi_{G_v^{ab},
q}(u)$). The argument is a contrapositive: if all three fail, a single nonzero $r \in K^D$
$\bigl(r := \sum_j \lambda_{(ab)j}\, r_j(q(ab))\bigr)$ is orthogonal to every supporting extensor in
the union~(6.45), which Lemma~2.1 forces to span all of $K^D$, so $r = 0$ --- a contradiction.

\begin{lemma}[The $D$ panel-support extensors of $k+2$ linearly-independent homogeneous vectors span $K^D$; KT eq.\ (6.45) via Lemma 2.1]
  \label{lem:case-III-claim612-extensor-span}
  \lean{CombinatorialRigidity.Molecular.span_omitTwoExtensor_eq_top}
  \leanok
  \uses{lem:extensor-independence, def:rigidity-matrix}
  Let $\bar p_1, \dots, \bar p_{k+2} \in K^{k+2}$ be linearly independent homogeneous vectors. The
  $\binom{k+2}{2} = D$ panel-support $k$-extensors $C(\bar p_i \bar p_j) = \bar p_i \vee \bar p_j$ ---
  one per pair, the codim-2 panel through each pair --- span $\bigwedge^k K^{k+2}$:
  \[
    \operatorname{span}\{\, \bar p_i \vee \bar p_j \mid 1 \le i < j \le k+2 \,\} = \bigwedge^k K^{k+2}
    \quad (\operatorname{finrank} D).
  \]
  At $d = 3$ this is the $k = 2$ case ($D = 6$, $\bigwedge^2 K^4$), the form used in the $d = 3$
  case below.
\end{lemma}
\begin{proof}
  \leanok
  By Lemma~2.1 (\cref{lem:extensor-independence}), applied at grade $e = k$ to any $k+2$ linearly
  independent (not necessarily affinely independent) vectors in $K^{k+2}$, the $D$ omit-two
  extensors $\bar p_i \vee \bar p_j$ are linearly independent in $\bigwedge^k K^{k+2}$. The
  strictly-increasing index pairs number $\binom{k+2}{2} = D = \operatorname{finrank} \bigwedge^k
  K^{k+2}$, so this linearly independent family is a basis, hence spans.
\end{proof}

\begin{lemma}[A functional annihilating a spanning set is zero]
  \label{lem:case-III-claim612-orthseq-vanish}
  \lean{CombinatorialRigidity.Molecular.eq_zero_of_annihilates_span_top}
  \leanok
  \uses{lem:case-III-claim612-extensor-span}
  If a functional $r \in (K^D)^* = (\bigwedge^k K^{k+2})^*$ on the screw space $\bigwedge^k
  K^{k+2}$ vanishes on every element of a set $S$ whose span is all of $K^D$, then $r = 0$. (The
  non-degeneracy of the dual pairing on $K^D$: a functional vanishing on a spanning set vanishes
  everywhere.)
\end{lemma}
\begin{proof}
  \leanok
  Two linear maps agreeing on a spanning set are equal; $r$ agrees with the zero functional on
  $S$, and $\operatorname{span} S = K^D$, so $r = 0$.
\end{proof}

The contrapositive needs the union of eq.~(6.45) to span $K^D$, which rests on two facts: first,
four \emph{affinely independent} points $p_1, \dots, p_4$ realizing the $\Pi(a)/\Pi(b)/\Pi(c)$
incidence pattern, so that Lemma~2.1 (\cref{lem:case-III-claim612-extensor-span}) forces their six
joins to span; and second, the \emph{point-join $\leftrightarrow$ panel-meet duality}, which turns
the hypothesis ``$r \perp C(L)$ for every line $L$ in $\Pi(a) \cup \Pi(b) \cup \Pi(c)$'' into ``$r$
annihilates each spanning join $p_i \vee p_j$''. The first is an existence fact --- one seed
realizing the projective construction suffices --- proved by the same non-root argument the
Case~I/II realizations use; the second is Grassmann--Cayley projective duality.

\begin{lemma}[The intersection of fewer than $d+1$ panels has a homogeneous representative; KT eq.\ (6.45), existence half]
  \label{lem:case-III-claim612-panel-point-exists}
  \lean{CombinatorialRigidity.Molecular.exists_ne_zero_dotProduct_eq_zero}
  \leanok
  \uses{def:panel-hinge-framework}
  Given $m < d + 1$ panel normals $n \colon \mathrm{Fin}\,m \to K^{d+1}$, the homogeneous incidence
  system $\langle \bar p, n_i \rangle = 0$ ($1 \le i \le m$) has a \emph{nonzero} solution
  $\bar p \in K^{d+1}$: the $m \times (d+1)$ coefficient matrix has strictly more columns than rows,
  so its null space is nontrivial. This is the existence half of KT eq.~(6.67) (p.~698): the
  intersection of any $j$ of the panels is a genuine $(d + 1 - j)$-dimensional affine subspace, and
  hence carries a homogeneous representative incident to all $j$.
\end{lemma}
\begin{proof}
  \leanok
  The incidence system is the linear map $K^{d+1} \to K^m$, $\bar p \mapsto (\langle \bar p, n_i
  \rangle)_i$, whose source dimension $d + 1$ strictly exceeds its target dimension $m$. By
  rank--nullity its kernel has dimension $\ge (d + 1) - m > 0$, so the kernel is nontrivial and
  contains a nonzero vector --- a homogeneous point orthogonal to every $n_i$.
\end{proof}

\begin{lemma}[A nonzero determinant polynomial in the seed yields an affinely-independent realization; KT eq.\ (6.45), the closure step]
  \label{lem:case-III-claim612-points-affineIndep-producer}
  \lean{CombinatorialRigidity.Molecular.exists_affineIndependent_of_det_polynomial_ne_zero}
  \leanok
  \uses{lem:affine-indep-iff}
  Suppose the four candidate points are built as functions $p \colon (\sigma \to K) \to
  \mathrm{Fin}\,4 \to K^3$ of a panel-coordinate seed $q \colon \sigma \to K$, and their
  affine-independence determinant --- the homogenization determinant
  $\det(\overline{p(q)}) = \det(\mathrm{homogenize} \circ p(q))$ of \cref{lem:affine-indep-iff} ---
  equals a \emph{nonzero} multivariate polynomial $P \in K[\sigma]$ evaluated at $q$. Then there is
  a seed $q$ at which the four points $p(q)$ are affinely independent.
\end{lemma}
\begin{proof}
  \leanok
  Over the infinite field $K$ a nonzero polynomial does not vanish identically, so $P$ has a
  non-root $q$, at which $\det(\overline{p(q)}) = P(q) \ne 0$; the determinant characterization
  (\cref{lem:affine-indep-iff}) turns this into affine independence of $p(q)$. This is the closure
  half of \cref{lem:case-III-claim612-points-affineIndep}: it isolates the existence-bearing content
  ($P \ne 0$) from the surrounding linear algebra, exactly as the dimension count
  (\cref{lem:case-III-claim612-panel-point-exists}) isolates the existence half of the argument.
\end{proof}

\begin{lemma}[The affine-independence determinant of polynomial-valued candidate points is a polynomial in the seed; KT eq.\ (6.45), determinant-polynomial bridge]
  \label{lem:case-III-claim612-detpoly}
  \lean{CombinatorialRigidity.Molecular.exists_detPolynomial_of_pointPolynomial}
  \leanok
  \uses{lem:affine-indep-iff}
  Suppose the four candidate points are built coordinate-by-coordinate as multivariate polynomials
  in the panel-coordinate seed: a family $pp \colon \mathrm{Fin}\,4 \to K^3[\sigma]$ of polynomial
  entries, with realized point $p(q)_i := \mathrm{eval}_q(pp_i)$. Then the affine-independence
  determinant of $p(q)$ --- the homogenization determinant $\det(\mathrm{homogenize} \circ p(q))$
  of \cref{lem:affine-indep-iff} --- is the evaluation at $q$ of a single polynomial
  $P \in K[\sigma]$: the determinant of the $(d+1) \times (d+1)$ polynomial matrix with rows
  $\overline{pp_i} := (pp_i, 1)$.
\end{lemma}
\begin{proof}
  \leanok
  Take $P$ to be the determinant of the polynomial matrix with rows $\overline{pp_i} = (pp_i, 1)$.
  Evaluation $\mathrm{eval}_q$ is a ring homomorphism, so it commutes with the determinant and with
  appending the constant coordinate $1$; hence $\mathrm{eval}_q(P)$ equals the determinant of the
  matrix of homogenized evaluated points $\mathrm{homogenize}(p(q)_i)$, which is exactly
  $\det(\overline{p(q)})$. This supplies the determinant hypothesis of
  \cref{lem:case-III-claim612-points-affineIndep-producer}, isolating the polynomial-in-the-seed
  structure of the determinant from the genericity-bearing fact that $P$ is nonzero.
\end{proof}

\begin{lemma}[The explicit good seed: four affinely-independent points realizing the $\Pi(a)/\Pi(b)/\Pi(c)$ incidence pattern; KT eq.\ (6.45), the $P \ne 0$ witness]
  \label{lem:case-III-claim612-points-affineIndep-witness}
  \lean{CombinatorialRigidity.Molecular.exists_affineIndependent_panel_incidence}
  \leanok
  \uses{lem:affine-indep-iff}
  There exist three panel normals $n \colon \mathrm{Fin}\,3 \to K^4$ in nonparallel position
  (linearly independent) and four \emph{affinely-independent} points $p \colon \mathrm{Fin}\,4 \to
  K^3$ realizing the triple-intersection incidence pattern of KT eq.~(6.45):
  $p_1 \in \Pi(a) \cap \Pi(b) \cap \Pi(c)$, $p_2 \in \Pi(a) \cap \Pi(b) \setminus \Pi(c)$,
  $p_3 \in \Pi(b) \cap \Pi(c) \setminus \Pi(a)$, $p_4 \in \Pi(c) \cap \Pi(a) \setminus \Pi(b)$, where
  panel incidence is $p_i \in \Pi(u) \Leftrightarrow \langle \bar{p_i}, n_u \rangle = 0$ on the
  homogenization $\bar{p_i}$.
\end{lemma}
\begin{proof}
  \leanok
  Take the coordinate-aligned seed: panel normals $n_a = e_0$, $n_b = e_1$, $n_c = e_2$ (the first
  three standard covectors of $K^4$, hence linearly independent, satisfying the nonparallel
  hypothesis), and the standard affine $3$-simplex $p = (0, e_3, e_1, e_2)$ of $K^3$ (the origin
  plus three axis points). The incidence pattern is immediate from the coordinates: each
  $\bar{p_i}$ vanishes against exactly the panel covectors whose coordinate it is zero at. The
  $4 \times 4$ homogenization determinant is $\pm 1 \ne 0$ (a direct cofactor expansion), so by
  \cref{lem:affine-indep-iff} the four points are affinely independent.
\end{proof}

\begin{lemma}[Four affinely-independent points with the $\Pi(a)/\Pi(b)/\Pi(c)$ incidence pattern; KT eq.\ (6.45) point choice]
  \label{lem:case-III-claim612-points-affineIndep}
  \lean{CombinatorialRigidity.Molecular.exists_affineIndependent_panel_incidence}
  \leanok
  \uses{def:rigidity-matrix, def:panel-hinge-framework,
        lem:case-III-claim612-panel-point-exists, lem:case-III-claim612-points-affineIndep-producer,
        lem:case-III-claim612-detpoly, lem:case-III-claim612-points-affineIndep-witness}
  At $d = 3$ there is a nonparallel arrangement of three panels --- three normals
  $n \colon \mathrm{Fin}\,3 \to K^4$ in nonparallel position (linearly independent), spanning
  panels $\Pi(a), \Pi(b), \Pi(c) \subset K^3$ --- together with four \emph{affinely-independent}
  points $p \colon \mathrm{Fin}\,4 \to K^3$ realizing the triple-intersection incidence pattern
  \[
    p_1 \in \Pi(a) \cap \Pi(b) \cap \Pi(c), \quad
    p_2 \in \Pi(a) \cap \Pi(b) \setminus \Pi(c), \quad
    p_3 \in \Pi(b) \cap \Pi(c) \setminus \Pi(a), \quad
    p_4 \in \Pi(c) \cap \Pi(a) \setminus \Pi(b),
  \]
  where membership is the panel-incidence relation $p_i \in \Pi(u) \Leftrightarrow
  \langle \bar{p_i}, n_u \rangle = 0$ on the homogenization $\bar{p_i}$. The four points span an
  affine $3$-simplex, so they are exactly the affinely-independent input
  \cref{lem:case-III-claim612-extensor-span} needs (their six joins $\bar{p_i} \vee \bar{p_j}$ then
  span $K^D$).
\end{lemma}
\begin{proof}
  \leanok
  This is an existence fact, like the Case~I/II realizations: it needs only that \emph{one} seed
  work, not that the inductive seed itself be generic. The explicit witness
  (\cref{lem:case-III-claim612-points-affineIndep-witness}) supplies it directly: the
  coordinate-aligned normals $n_a = e_0$, $n_b = e_1$, $n_c = e_2$ are linearly independent in
  $K^4$, hence nonparallel, and the standard affine $3$-simplex $p = (0, e_3, e_1, e_2)$ of $K^3$
  has homogenization determinant $\pm 1 \ne 0$ (\cref{lem:affine-indep-iff}), with the incidence
  pattern read straight off the coordinates.

  The general argument --- constructing the four points as polynomials in an arbitrary normal
  arrangement $n$ and showing the resulting determinant polynomial is nonzero
  (\cref{lem:case-III-claim612-points-affineIndep-producer} and
  \cref{lem:case-III-claim612-detpoly}, together
  with the dimension count \cref{lem:case-III-claim612-panel-point-exists}) --- would also work, by
  the equivalence between a nonzero determinant polynomial and one seed at which the constructed
  points are affinely independent; the explicit witness above makes it unnecessary by fixing the
  arrangement outright.

  The panel normals are existentially quantified because the Case~III realization is itself being
  constructed: the candidate placements $p_1, p_2, p_3$
  (\cref{lem:case-II-realization-placement}, \cref{lem:case-III-claim612-p2-placement}, and
  \cref{lem:case-III-claim612-p3-placement}) realize $G_v^{ab}$ at a panel-coordinate seed $q$ that
  \cref{lem:case-III-claim612} is free to choose, so the points and the panels they are incident to
  come from one chosen realization, not from a fixed external normal arrangement.
\end{proof}

The witness of \cref{lem:case-III-claim612-points-affineIndep-witness} fixes the three panel
normals to the coordinate frame $e_0, e_1, e_2$, but \cref{lem:case-II-realization} reads its
three panel normals $n_a, n_b, n_c$ off the split-leg realization $q$ it is handed --- nonparallel,
but otherwise arbitrary --- and needs a witness relative to \emph{those} normals. The following
lemma supplies one directly, in homogeneous coordinates rather than by de-homogenizing to affine
points: the two later uses --- the span \cref{lem:case-III-claim612-extensor-span} (via Lemma~2.1,
which needs only linear independence of the four homogeneous vectors, not affine independence of
the underlying points) and the point-join/panel-meet transfer
\cref{lem:case-III-claim612-line-in-panel-union} (which reads the joins as bare vectors with
$\langle \cdot, n_u \rangle = 0$ incidence) --- both need only a linearly independent family in
$K^4$.

\begin{lemma}[The witness points, parameterized by the actual split-leg panel normals; KT eq.\ (6.45)]
  \label{lem:case-III-claim612-points-affineIndep-real-normals}
  \lean{CombinatorialRigidity.Molecular.exists_homogeneousIncidence_of_normals}
  \leanok
  \uses{lem:case-III-claim612-panel-point-exists, lem:affine-indep-iff}
  For any linearly independent panel normals $n \colon \mathrm{Fin}\,3 \to K^4$, there exist four
  linearly independent homogeneous vectors $\bar p \colon \mathrm{Fin}\,4 \to K^4$ realizing the
  same triple-intersection incidence pattern relative to $n$: $\langle \bar p_1, n_u \rangle = 0$
  for every $u$, and each $\bar p_{i+1}$ orthogonal to two of the $n_u$ and nonorthogonal to the
  third.
\end{lemma}
\begin{proof}
  \leanok
  The pairing map $x \mapsto (u \mapsto \langle x, n_u\rangle)$ is the $\mathrm{mulVecLin}$ of the
  row matrix whose rows are the $n_u$; its rows are linearly independent, so the matrix has full row
  rank $3$ and the map is \emph{surjective}. Hence for any prescribed pairing target $t \in K^3$
  there is a preimage $x$ with $\langle x, n_u\rangle = t_u$. Take $\bar p_2, \bar p_3, \bar p_4$ as
  the preimages of $(0,0,1), (1,0,0), (0,1,0)$ (giving exactly the off-one-panel incidence), and
  $\bar p_1$ as the nonzero common-perp vector of all three normals
  (\cref{lem:case-III-claim612-panel-point-exists}, $3 < 4$). Linear independence of the four
  $\bar p_i$ is the triangular argument on the incidence matrix: pairing
  $\sum_i g_i \bar p_i = 0$ against $n_0$ kills all but $\bar p_3 \Rightarrow g_3 = 0$; against $n_1
  \Rightarrow g_4 = 0$; against $n_2 \Rightarrow g_2 = 0$; then $g_1 \bar p_1 = 0$ with
  $\bar p_1 \ne 0 \Rightarrow g_1 = 0$. No genericity of $n$ beyond linear independence is used.
\end{proof}

The point-join $\leftrightarrow$ panel-meet duality itself --- that for two points
$p_i, p_j$ on a line $L$ lying in a panel $\Pi(u)$, the point-join $\bar{p_i} \vee \bar{p_j}$ and
the panel-meet supporting extensor of $L$ are scalar multiples of one another, so a functional
annihilating every panel-meet of a line in $\Pi(a) \cup \Pi(b) \cup \Pi(c)$ also annihilates every
such spanning join --- is proved once, at the level of the Grassmann--Cayley meet, as
\cref{lem:case-III-claim612-line-in-panel-union}.

\begin{lemma}[Block full rank $\Leftrightarrow$ $\hat r \notin (\operatorname{span} C)^\perp$; KT (6.42) row-space criterion]
  \label{lem:case-III-claim612-block-iff-perp}
  \lean{CombinatorialRigidity.Molecular.BodyHingeFramework.linearIndependent_sumElim_candidateRow_iff,
        CombinatorialRigidity.Molecular.BodyHingeFramework.mem_hingeRowBlock_iff}
  \leanok
  \uses{def:rigidity-matrix, def:hinge-row-block, lem:case-II-placement-new-rows}
  Let the $D - 1$ rows of a candidate's hinge block span $(\operatorname{span} C)^\perp = r(p(e))$
  (the $D-1$-dimensional orthogonal complement of the line $\operatorname{span} C$ supporting the
  candidate's $(d-2)$-affine subspace $L$; \cref{lem:case-II-placement-new-rows}), and let
  $\hat r := \sum_j \lambda_{(ab)j}\, r_j(q(ab))$ be the candidate row. The augmented $D$-functional
  family ($D - 1$ block rows plus $\hat r$) is linearly independent iff
  \[
    \hat r \notin (\operatorname{span} C)^\perp \quad\Longleftrightarrow\quad \hat r(C) \ne 0,
  \]
  i.e.\ the candidate's top-left $D \times D$ block has full rank iff $\hat r$ is not orthogonal to
  the supporting extensor $C$.
\end{lemma}
\begin{proof}
  \leanok
  The $D - 1$ block rows already span the hyperplane $(\operatorname{span} C)^\perp$; appending
  $\hat r$ raises the dimension to $D$ iff $\hat r$ is fresh, i.e.\ $\hat r \notin
  (\operatorname{span} C)^\perp$ (else the family stays inside the $(D-1)$-dimensional span).
  Freshness off the hinge-row block is exactly $\hat r(C) \ne 0$, by the dual-annihilator
  correspondence $\hat r \in (\operatorname{span} C)^\perp = r(p(e)) \Leftrightarrow \hat r(C) = 0$
  (the annihilator evaluated at the spanning vector $C$).
\end{proof}

\begin{lemma}[The common vector $r$ is nonzero; KT (6.42)]
  \label{lem:case-III-claim612-r-nonzero}
  \lean{CombinatorialRigidity.Molecular.BodyHingeFramework.candidateRow_ne_zero}
  \leanok
  \uses{lem:case-III-redundant-decomposition, lem:case-II-placement-new-rows}
  The common vector of the disjunction, $r := \sum_{1 \le j \le 5} \lambda_{(ab)j}\, r_j(q(ab))
  \in K^D$, is nonzero. (The $r_j(q(ab))$ are the $D - 1$ linearly independent rows of the $ab$-hinge
  block, and the coefficient $\lambda_{(ab)i^*} = 1$ from the redundant-row decomposition of
  eq.~(6.25); a combination with a unit coefficient on a member of a linearly independent family is
  nonzero.)
\end{lemma}
\begin{proof}
  \leanok
  The $r_j(q(ab))$ ($1 \le j \le 5$, i.e.\ $D - 1$ of them) are linearly independent
  (\cref{lem:case-II-placement-new-rows}); by \cref{lem:case-III-redundant-decomposition} the
  coefficient $\lambda_{(ab)i^*} = 1$ (eq.~(6.25)). A linear combination of a linearly independent
  family with a nonzero coefficient on one member is nonzero.
\end{proof}

\begin{lemma}[The $p_2$ placement: the symmetric $va \leftrightarrow vb$ candidate; KT eqs.\ (6.19)/(6.30)]
  \label{lem:case-III-claim612-p2-placement}
  \lean{CombinatorialRigidity.Molecular.BodyHingeFramework.linearIndependent_sum_p2_candidateRow}
  \leanok
  \uses{lem:case-II-realization-placement, lem:case-III-columnop,
        lem:case-III-claim612-block-iff-perp}
  The second candidate $p_2$: split off at $v$ along $vb$ (rather than $va$), the symmetric image of
  $p_1$ under $a \leftrightarrow b$. It reuses the eq.~(6.12) degenerate placement
  (\cref{lem:case-II-realization-placement}) with the roles of $a, b$ exchanged, and the column
  operation at $(v, b)$ (\cref{lem:case-III-columnop}) in place of $(v, a)$. By the same
  candidate-completion assembly (\cref{lem:case-III-candidate-row}, applied at the swapped roles),
  \emph{if} its top-left $D \times D$ block $M_2$ (eq.~(6.30)) is full rank --- equivalently
  $r \notin (\operatorname{span} C(L'))^\perp$, i.e.\ $r(C(L')) \ne 0$ for the line $L' \subset
  \Pi(b)$ (the row-space criterion \cref{lem:case-III-claim612-block-iff-perp}) --- then $R(G, p_2)$
  attains the full $D(|V|-1)$-size independent family. The same $\lambda_{(ab)j}$ and index $i^*$
  appear in $M_1$ and $M_2$ (they are determined by $(G_v^{ab}, q)$, independent of $p_1, p_2$).
\end{lemma}
\begin{proof}
  \leanok
  Verbatim symmetric to the $p_1$ construction under $a \leftrightarrow b$. The candidate-completion
  assembly (\cref{lem:case-III-candidate-row}) applies abstractly to the splitting body, its two
  endpoints, and the candidate functional, so taking the splitting body to be $v$ and the column
  operation to be $\Phi = \mathrm{col}_b \mathrel{+}= \mathrm{col}_v$ for the edge $vb$, in place of
  $\mathrm{col}_a \mathrel{+}= \mathrm{col}_v$ for $va$, is the same lemma. Its one hypothesis ---
  the operated, $v$-pinned top-left block being full rank --- is supplied by the row-space criterion
  (\cref{lem:case-III-claim612-block-iff-perp}) at the $vb$-hinge once the operated candidate row is
  identified with $\rho$ itself (the column operation makes it pure $v$-column with value
  $\rho(S_v)$), which holds iff $\rho(C) \ne 0$. The redundant-row combination is unchanged: it
  lives in $R(G_v^{ab}, q)$, common to both candidates.
\end{proof}

\begin{lemma}[Link correspondence under the $\rho = (a\,v)$ relabel; KT eq.\ (6.31)]
  \label{lem:splitOff-isLink-relabel}
  \lean{Graph.splitOff_isLink_relabel}
  \leanok
  \uses{def:graph-operations}
  Let $v, a, b, c$ and $e_a, e_b, e_c$ be the vertices and edges of the chain $b$--$v$--$a$--$c$
  introduced above: $e_a$ links $v$ to $a$, $e_b$ links $v$ to $b$, $e_c$ links $a$ to $c$, every
  $G$-edge at $v$ is $e_a$ or $e_b$, and every $G$-edge at $a$ is $e_a$ or $e_c$. Fix fresh edges
  $e_0, e_1 \notin E(G)$ with $e_1 \ne e_0$, and write $\rho = (a\,v)$ for the transposition of $a$
  and $v$, and $\sigma = (e_b\,e_0)(e_1\,e_c)$ for the product of the two edge transpositions.
  Then an edge $e$ links $x$ to $y$ in the $a$-split $G_a^{vc}$ if and only if $\sigma e$ links
  $\rho x$ to $\rho y$ in the $v$-split $G_v^{ab}$. This is KT eq.~(6.31): the graph isomorphism $G_a^{vc} \cong
  G_v^{ab}$ sends $e_b \mapsto e_0$ (the surviving $v$--$b$ edge becomes the fresh $a$--$b$ edge)
  and $e_1 \mapsto e_c$ (the fresh $v$--$c$ edge becomes the surviving $a$--$c$ edge).
\end{lemma}
\begin{proof}
  \leanok
  Expand the link relation of each split-off graph and match cases using the degree-2 closures at
  $v$ and $a$. Surviving edges $e \notin \{e_b, e_1\}$ have both endpoints away from $v$ and $a$
  (by the two closures), so $\sigma$ and $\rho$ fix them. The $e_b$-case maps to the fresh
  $e_0$-edge (endpoints $v \mapsto a$, $b \mapsto b$); the $e_1$-case maps to the surviving
  $e_c$-edge (endpoints $v \mapsto a$, $c \mapsto c$).
\end{proof}

\begin{lemma}[Fixed-seed realization transport across the $\rho = (a\,v)$ relabel; KT eq.\ (6.31)]
  \label{lem:splitOff-ofNormals-relabel}
  \lean{CombinatorialRigidity.Molecular.PanelHingeFramework.ofNormals_relabel,
        CombinatorialRigidity.Molecular.PanelHingeFramework.hasGenericFullRankRealization_of_splitOff_relabel}
  \leanok
  \uses{lem:splitOff-isLink-relabel}
  Because the eq.~(6.44) argument below tests the same vector $\hat r$ against all three candidate
  panels, the $a$-split realization must share the $v$-split's seed via $\rho$, rather than being
  reconstructed independently (an independent realization would carry a different seed, hence
  different coefficients $\lambda$, collapsing the trichotomy); so the transport below is stated at
  the level of the free-normal panel framework itself, not merely as an implication between
  existence claims.

  The induction hypothesis realizes the $v$-split $G_v^{ab}$ as a free-normal panel framework with
  some endpoint selector $\mathrm{ends}_0$ and seed $q_0$, satisfying three conjuncts: general
  position, rigidity on $V(G) \setminus \{v\}$, and link-recording (every link of $G_v^{ab}$ is
  among the selector's pairs). Under the hypotheses of
  \cref{lem:splitOff-isLink-relabel}, the $a$-split $G_a^{vc}$ satisfies the same three conjuncts
  for the free-normal panel framework with the transported selector $\mathrm{ends}_{\sigma\rho}(e) :=
  (\rho\,(\mathrm{ends}_0(\sigma e))_1,\, \rho\,(\mathrm{ends}_0(\sigma e))_2)$ and the transported
  seed $q_\rho(x, i) := q_0(\rho x, i)$, where $\rho = (a\,v)$ and $\sigma = (e_b\,e_0)(e_1\,e_c)$
  as in \cref{lem:splitOff-isLink-relabel}. The bare existential form ($G_v^{ab}$ has a generic full-rank
  realization $\Rightarrow$ so does $G_a^{vc}$, at the seed $q_0 \circ \rho$) is recorded as a
  short corollary.
\end{lemma}
\begin{proof}
  \leanok
  \uses{lem:splitOff-isLink-relabel}
  Each of the three conjuncts transports along the vertex-swap $\rho$ and edge-permutation $\sigma$
  (an involution, since its two transpositions have disjoint supports). General position: the
  $\rho$-reindexed seed $q_\rho$ is in general position because $\rho$ is injective and $q_0$ is.
  Rigidity: any motion $S$ of the $a$-split framework pulls back to a motion $S \circ \rho$ of the
  $v$-split framework, using \cref{lem:splitOff-isLink-relabel} to move each $a$-split link to a
  $v$-split link and matching the two frameworks' supporting extensors; the $v$-split's rigidity on
  $V(G) \setminus \{v\}$ then forces $S \circ \rho$ constant, hence $S$ constant on
  $V(G) \setminus \{a\}$, since $\rho$ maps $V(G) \setminus \{a\}$ to $V(G) \setminus \{v\}$
  bijectively ($\rho \circ \rho = \mathrm{id}$). Link-recording: \cref{lem:splitOff-isLink-relabel}
  translates each $a$-split link forward to a $v$-split link whose endpoints $\mathrm{ends}_0$
  records, transported back through $\rho$ to $\mathrm{ends}_{\sigma\rho}$.
\end{proof}

\begin{lemma}[Row-space correspondence under the $\rho = (a\,v)$ relabel; KT eqs.\ (6.31)/(6.44)]
  \label{lem:splitOff-rigidityRows-relabel}
  \lean{CombinatorialRigidity.Molecular.PanelHingeFramework.rigidityRows_ofNormals_relabel}
  \leanok
  \uses{lem:splitOff-isLink-relabel, lem:splitOff-ofNormals-relabel}
  This is the row-space correspondence the $M_3$ candidate row (\cref{lem:case-III-claim612-p3-placement})
  needs. Write $\Psi_\rho$ for the relabeling map on screw assignments, $\Psi_\rho(S) := S \circ
  \rho$, and $\Psi_\rho^*$ for its dual on functionals, $\Psi_\rho^*(\phi) := \phi \circ \Psi_\rho$.
  Under $\Psi_\rho^*$, the rigidity rows of the relabelled framework on $G_a^{vc}$ (endpoint
  selector $\mathrm{ends}_{\sigma\rho}$, seed $q_\rho$, as in \cref{lem:splitOff-ofNormals-relabel})
  are exactly the image of the $v$-split framework's rigidity rows: the $a$-split row at endpoints
  $u, w$ with functional $r$ --- $S \mapsto r(S_u - S_w)$ --- is $\Psi_\rho^*$ of the $v$-split row
  at endpoints $\rho u, \rho w$ with the same $r$.
\end{lemma}
\begin{proof}
  \leanok
  \uses{lem:splitOff-isLink-relabel}
  A direct computation using $\rho \circ \rho = \mathrm{id}$ shows $\Psi_\rho^*$ sends the
  $v$-split row at $(\rho u, \rho w, r)$ to the $a$-split row at $(u, w, r)$: evaluated at a screw
  assignment $S$, both sides read $r(S_u - S_w)$. \Cref{lem:splitOff-isLink-relabel} puts the two
  graphs' links in $\rho$-correspondence, and the hinge-row blocks at corresponding edges coincide
  (the same supporting extensor, via the same seed reindex as in
  \cref{lem:splitOff-ofNormals-relabel}, since $\sigma$ is an involution). Hence the two
  rigidity-row sets correspond exactly under $\Psi_\rho^*$, by mutual inclusion.
\end{proof}

\begin{lemma}[The $p_3$ placement: the third candidate via $G_v^{ab} \cong G_a^{vc}$; KT eqs.\ (6.31)--(6.41)]
  \label{lem:case-III-claim612-p3-placement}
  \lean{CombinatorialRigidity.Molecular.BodyHingeFramework.linearIndependent_sum_p3_candidateRow}
  \leanok
  \uses{lem:case-II-realization-placement, lem:case-III-columnop,
        lem:case-III-claim612-block-iff-perp, lem:case-III-claim612-eq644}
  The third candidate $p_3$, available because $a$ too is a degree-2 vertex: split off at $a$ along
  $vc$. The resulting graph $G_a^{vc}$ is isomorphic to $G_v^{ab}$ by $\rho \colon V \setminus \{a\}
  \to V \setminus \{v\}$, $\rho(v) = a$, $\rho(u) = u$ otherwise; the induced placement
  $q_\rho$ satisfies $p_3(va) = q(ac)$, $p_3(vb) = q(ab)$, $p_3(ac) = L''$, $p_3(e) = q(e)$ otherwise
  (eq.~(6.33)). The column operation at $(a, c)$ ($\mathrm{col}_a \mathrel{+}= \mathrm{col}_c$,
  eq.~(6.35)) block-triangularizes $R(G, p_3)$ to eq.~(6.41), whose top-left $D \times D$ block is
  $M_3$ with rows $r(p_3(ac)) = r(L'')$ and $\sum_j \lambda_{(ac)j}\, r_j(q(ac))$, full rank iff
  $r \notin (\operatorname{span} C(L''))^\perp$ for $L'' \subset \Pi(c)$ --- using eq.~(6.44)
  (\cref{lem:case-III-claim612-eq644}) to identify the $M_3$ row with the same $r$.
\end{lemma}
\begin{proof}
  \leanok
  The isomorphism $\rho$ identifies $(G_a^{vc}, q_\rho)$ with $(G_v^{ab}, q)$ as panel-hinge
  frameworks (eq.~(6.31)), so $\Pi_{G_a^{vc}, q_\rho}(u) = \Pi(u)$ for $u \ne v, a$ and
  $\Pi_{G_a^{vc}, q_\rho}(v) = \Pi(a)$. The argument is then verbatim symmetric to the $p_2$
  construction (\cref{lem:case-III-claim612-p2-placement}) under the role change $v, b \mapsto a,
  c$: the candidate-completion assembly (\cref{lem:case-III-candidate-row}) applies abstractly to
  the splitting body, its two endpoints, and the candidate functional, so taking the splitting
  body to be $a$ and the column operation to be $\Phi = \mathrm{col}_a \mathrel{+}= \mathrm{col}_c$
  is the same lemma; no separate isomorphism step is needed. Its one hypothesis --- the operated,
  $a$-pinned top-left block being full rank --- is supplied by the row-space criterion
  (\cref{lem:case-III-claim612-block-iff-perp}) at the $ac$-hinge once the operated candidate row is
  identified with $\rho$ itself (the column operation makes it pure $a$-column with value
  $\rho(S_a)$), which holds iff $\rho(C) \ne 0$. The column operation at $(a,c)$ (eqs.~(6.35)--(6.36))
  and the substitutions $p_3(va) = q(ac)$, $p_3(vb) = q(ab)$ (eq.~(6.34)) reduce $R(G, p_3)$ to
  eq.~(6.41) by the same row operations as for $R(G, p_1)$ (eq.~(6.40)). The eq.~(6.44) identity
  (\cref{lem:case-III-claim612-eq644}) gives $M_3$'s row as $-\sum_j \lambda_{(ac)j}\, r_j(q(ac)) =
  r$ (up to sign), so the full-rank criterion for $M_3$ reads off the same common vector $r$.
\end{proof}

\begin{lemma}[Eq.\ (6.44): the $M_3$ row equals $-r$; KT eq.\ (6.44)]
  \label{lem:case-III-claim612-eq644}
  \lean{CombinatorialRigidity.Molecular.BodyHingeFramework.candidateRow_ac_eq_neg,
        CombinatorialRigidity.Molecular.BodyHingeFramework.hingeRow_comp_single_tail,
        CombinatorialRigidity.Molecular.BodyHingeFramework.hingeRow_comp_single_off,
        CombinatorialRigidity.Molecular.BodyHingeFramework.acolumn_mem_hingeRowBlock_of_span_rigidityRows,
        CombinatorialRigidity.Molecular.BodyHingeFramework.hingeRow_acolumn_mem_span_rigidityRows}
  \leanok
  \uses{lem:case-III-acolumn-zero, lem:case-III-redundant-decomposition}
  In $G_v^{ab}$ only $ab$ and $ac$ are incident to the body $a$, so the $a$-column block of the
  eq.~(6.24) vanishing combination (\cref{lem:case-III-acolumn-zero}, eq.~(6.43)) reads
  \[
    0 = \sum_{1 \le j \le 5} \lambda_{(ab)j}\, r_j(q(ab)) + \sum_{1 \le j \le 5} \lambda_{(ac)j}\,
        r_j(q(ac)),
  \]
  which, with $r = \sum_j \lambda_{(ab)j}\, r_j(q(ab))$, gives KT's eq.~(6.44):
  \[
    r = -\sum_{1 \le j \le 5} \lambda_{(ac)j}\, r_j(q(ac)).
  \]
  Hence the $M_3$ candidate row $\sum_j \lambda_{(ac)j}\, r_j(q(ac))$ equals $-r$, so the third
  candidate's full-rank criterion reads off the \emph{same} common vector $r$ as $M_1, M_2$.
\end{lemma}
\begin{proof}
  \leanok
  Eq.~(6.43) (\cref{lem:case-III-acolumn-zero}) is the $a$-column restriction of the eq.~(6.24)
  vanishing combination; since only the $ab$- and $ac$-edges are incident to $a$ in $G_v^{ab}$,
  every other row's $a$-column entry vanishes by the rigidity-matrix definition, and only the
  $ab$- and $ac$-terms survive the restriction. Rearranging $\sum_j \lambda_{(ab)j}\, r_j(q(ab)) +
  \sum_j \lambda_{(ac)j}\, r_j(q(ac)) = 0$ for $r$ is eq.~(6.44).
\end{proof}

\begin{lemma}[Claim 6.12: for some line choice the candidate block has full rank; KT Claim 6.12, eqs.\ (6.30)--(6.45)]
  \label{lem:case-III-claim612}
  \lean{CombinatorialRigidity.Molecular.BodyHingeFramework.case_III_claim612_gen}
  \leanok
  \uses{lem:case-III-claim612-points-affineIndep-real-normals, lem:case-III-claim612-r-nonzero}
  At $d = 3$, let $r \ne 0$ be the common candidate row and let $\bar p_1, \bar p_2, \bar p_3,
  \bar p_4 \in K^4$ be four linearly independent homogeneous vectors
  (\cref{lem:case-III-claim612-points-affineIndep-real-normals}). Then for some pair $(i, j)$ the
  join $r(\bar p_i \vee \bar p_j) \ne 0$. Since a candidate's hinge support is the panel-meet of
  its (freely chosen) line, and the panel-meet of a line through two of the four points is exactly
  their join $p_i \vee p_j$, this is KT's Claim~6.12: for some choice of $(d-2)$-affine line, the
  corresponding candidate's top-left $D \times D$ block (one of $M_1, M_2, M_3$ of eqs.~(6.29),
  (6.30), (6.41)) has full rank, equivalently $R(G, p_i)$ attains the desired rank $D(|V| - 1)$ ---
  discharging the one conditional the candidate-completion above leaves open.
\end{lemma}
\begin{proof}
  \leanok
  \uses{lem:case-III-claim612-extensor-span, lem:case-III-claim612-orthseq-vanish}
  Contrapositive of the existential claim. Suppose $r$ annihilated every one of the six joins
  $\bar p_i \vee \bar p_j$. (This is precisely KT's eq.~(6.45) hypothesis ``for every choice of
  lines $L \subset \Pi(a)$, $L' \subset \Pi(b)$, $L'' \subset \Pi(c)$'', once each join is
  identified with the panel-meet $C(L)$ of the line through that pair via the point-join
  $\leftrightarrow$ panel-meet duality (\cref{lem:case-III-claim612-line-in-panel-union}, with
  eq.~(6.44) placing all three candidate blocks' failure conditions on the same $r$); the argument
  works directly with the joins.) But the six joins of four linearly independent homogeneous
  vectors span $K^D$ (\cref{lem:case-III-claim612-extensor-span}, Lemma~2.1), so $r$ would
  annihilate a spanning set and hence be $0$ (\cref{lem:case-III-claim612-orthseq-vanish}),
  contradicting $r \ne 0$ (\cref{lem:case-III-claim612-r-nonzero}). Therefore
  $r(\bar p_i \vee \bar p_j) \ne 0$ for some pair.

  The claim is genuinely an existential over a free choice of lines, not a disjunction over three
  fixed supports: three $2$-extensors $C_1, C_2, C_3$ span at most $3$ of the $6$ dimensions of
  $\bigwedge^2 K^4$, so $r \perp C_1, C_2, C_3$ alone cannot force $r = 0$; only the full six-join
  span of eq.~(6.45) does. The realization is then built to exploit the witness pair directly: its
  candidate placement takes the hinge line to \emph{be} the witness join's line
  $L = \bar p_i \bar p_j$.
\end{proof}

Assembling the pieces gives the three-dimensional case directly. Claim~6.12
(\cref{lem:case-III-claim612}) names one of the three candidates whose top-left block is
full rank; the candidate-completion (\cref{lem:case-III-candidate-row}) then realizes $G$ at
that candidate's degenerate placement. This is the concrete counterpart of the general chain
dispatch (\cref{lem:case-III-chain-dispatch}, below): where the general argument transports a
single redundancy along a chain to reach the discriminating panel, the $d = 3$ case
dispatches it to one of three fixed candidates.

\begin{lemma}[The $d = 3$ three-candidate dispatch; KT Lemma 6.10, eqs.\ (6.24)--(6.44)]
  \label{lem:case-III-candidate-dispatch-d3}
  \lean{CombinatorialRigidity.Molecular.PanelHingeFramework.case_III_candidate_dispatch}
  \leanok
  \uses{def:panel-hinge-framework, def:rank-hypothesis, lem:case-III-nested-rank-lower}
  Fix $d = 3$ ($k = 2$, $D = 6$). Let $G$ be a simple $0$-dof-graph with an adjacent pair of
  degree-$2$ vertices --- $v$ with neighbours $a, b$ (incident edges $e_a \ne e_b$) and $a$
  with neighbours $v, c$ (incident edges $e_a \ne e_c$) --- and a fresh edge $e_0 \notin
  E(G)$. Suppose the split-off $G_v^{ab}$ is $0$-dof and carries a generic full-rank
  realization (the induction hypothesis at the smaller graph), and the eq.~(6.22) nested
  rank bound holds at the vertex-removal $G - v$
  (\cref{lem:case-III-nested-rank-lower}). Then $G$ carries a generic full-rank realization
  at rank $D(|V| - 1)$.
\end{lemma}
\begin{proof}
  \leanok
  \uses{lem:case-III-claim-6-11, lem:case-III-candidate-row, lem:case-III-claim612,
        lem:case-III-claim612-p2-placement, lem:case-III-claim612-p3-placement,
        lem:case-II-realization-placement, lem:genericity-device}
  From the generic realization of the split-off $G_v^{ab}$ at a seed $q$, Claim~6.11
  (\cref{lem:case-III-claim-6-11}) and the eq.~(6.22) nested rank bound
  (\cref{lem:case-III-nested-rank-lower}) produce the shared redundancy: a nonzero
  functional $\rho$ annihilating the $e_0$-extensor $C(q(ab))$, whose $ab$-hinge row lies in
  $\operatorname{span}(R(G_v, q)\text{-rows})$, together with the $D(|V_v| - 1)$ independent
  bottom rows the candidate-completion needs (\cref{lem:case-III-candidate-row}). Reading the
  three candidate panels $\Pi(a), \Pi(b), \Pi(c)$ off $q$, Claim~6.12
  (\cref{lem:case-III-claim612}) selects a panel and a transversal direction at which $\rho$
  is not annihilated. A case split over the three panels realizes $G$ at the matched
  candidate: the $a$-side candidate is the eq.~(6.12) degenerate placement
  (\cref{lem:case-II-realization-placement}) completed by \cref{lem:case-III-candidate-row};
  the $b$-side candidate is its symmetric image under $a \leftrightarrow b$
  (\cref{lem:case-III-claim612-p2-placement}); and the $c$-side candidate is the third
  placement obtained through the relabel $G_v^{ab} \cong G_a^{vc}$
  (\cref{lem:case-III-claim612-p3-placement}). Each attains rank $D(|V| - 1)$ at a degenerate
  placement, and the genericity device (\cref{lem:genericity-device}) re-realizes that rank
  at a general-position seed.
\end{proof}
\begin{fmlnote}
  The eq.~(6.22) nested rank bound rides here as an explicit hypothesis, discharged by
  \cref{lem:case-III-nested-rank-lower}; the general \cref{lem:case-III} derives the same
  bound internally from its induction hypothesis. This is a worked concrete case: the general
  induction runs the chain dispatch (\cref{lem:case-III-chain-dispatch}) rather than this
  three-candidate specialization, so the pinned declaration and its supporting placements
  sit off the proof's live path. They are retained deliberately, in the Lean as in this
  chapter, not as a superseded draft: the fixed three-candidate dispatch is simpler than
  the chain transport, it is the form KT themselves give for $d = 3$ (\S6.4.1), and it is
  kept as the accessible worked instance of the general argument.
\end{fmlnote}

\medskip\noindent\emph{The triangle floor.} The Case-III recursion needs a floor
of its own, and finding it is a point where the formalization must do visible
work that the paper covers in a phrase. KT open Lemma~6.13 by Lemma~4.6: $G$ is
either a short cycle (handled by the Lemma~5.4 cycle family,
\cref{lem:cycle-realization}) or carries a length-$d$ chain. The formalized
reduction runs this dichotomy directly (\cref{lem:chain-data-extract}); at
$d \le 3$ its cycle disjunct is mathematically vacuous (a cycle witness would
force $|V(G)| \le d \le 3$, contradicting the $|V(G)| \ge 4$ floor the chain
disjunct runs under), so in practice the split recursion always descends via
the chain disjunct until $|V(G)| = 3$. There it genuinely cannot recurse:
splitting off a degree-2 vertex of a triangle leaves the \emph{double edge} on
two bodies --- a non-simple graph --- so the conditioned induction hypothesis
(``general position \emph{if simple}'') supplies only the bare conjunct for
it, and with it the candidate placement loses its transversality input (the
linear independence of the two split-leg normals). A triangle is a minimal
$0$-dof-graph that cannot recurse further, so its realization is produced
\emph{directly} --- and a triangle is precisely the $m = 3$ cycle. The floor is
therefore realized as the $m = 3$ instance of the Lemma~5.4 cycle family: pin
the vertex set and find the third edge by edge counting
(\cref{lem:triangle-third-edge}), package $G$ as three-cycle data, and apply
\cref{lem:cycle-realization}. This is exactly the source's shape: KT state
Lemma~5.4 at $3 \le |V| \le D$ and never treat the triangle as a separate
case. A standalone three-vertex assembly (\cref{lem:triangle-realization},
realizing three mutually transversal panels by a cyclically consistent
standard-basis seed (\cref{lem:triangle-normals}), concluding rigidity from the
three-extensor telescoping independence (\cref{lem:theorem-55-triangle}, the
three-body sibling of the two-body base case \cref{lem:theorem-55-base}), and
upgrading by the genericity device) is also given, as the accessible worked
instance of the general cyclic construction, off the recursion's live path.
The one Lemma~5.4 node (\cref{lem:cycle-realization}) thus serves both the
triangle floor at $m = 3$ and the $|V(G)| \ge 4$ chain dichotomy's cycle
disjunct at $m \ge 4$ (the latter vacuous at $d \le 3$).

\begin{lemma}[$|V|=3$ triangle base: vertex and edge pin; KT \S6.4]
  \label{lem:triangle-third-edge}
  \lean{Graph.exists_isLink_of_isMinimalKDof_card_three}
  \leanok
  \uses{def:k-dof}
  Let $G$ be a simple minimal $0$-dof graph on exactly three vertices, with edges $e_a : v$--$a$
  and $e_b : v$--$b$ and $e_a \ne e_b$. Then $a \ne b$, $V(G) = \{v,a,b\}$, and there exists an
  edge $f : a$--$b$.
\end{lemma}
\begin{proof}
  \leanok
  \uses{thm:def-eq-corank}
  (1) $a \ne b$: if $a = b$ then $e_a, e_b$ are parallel, contradicting simplicity.
  (2) Vertex pin: $\{v,a,b\} \subseteq V(G)$ from the edge hypotheses; $|\{v,a,b\}| = 3 = |V(G)|$
  forces equality.
  (3) Third edge: the rank formula $\operatorname{rank} M(\tilde G) = D(|V|-1) = 2D$
  (\cref{thm:def-eq-corank}, at deficiency $0$) and the rank bound
  $\operatorname{rank} M(\tilde G) \le (D-1)|E(G)|$ (rank is at most the ground-set size) force
  $|E(G)| \ge 3$, so a third edge $f \notin \{e_a, e_b\}$ exists. Its endpoints lie in
  $\{v,a,b\}$; simplicity --- no two edges share both endpoints --- eliminates every case but
  $f : a$--$b$.
\end{proof}

\begin{lemma}[$|V|=3$ triangle base: generic realization; KT \S6.4]
  \label{lem:triangle-realization}
  \lean{CombinatorialRigidity.Molecular.PanelHingeFramework.hasGenericFullRankRealization_of_triangle}
  \leanok
  \uses{lem:triangle-third-edge, lem:theorem-55-triangle, lem:triangle-normals}
  Let $G$ be a simple minimal $0$-dof graph on exactly three vertices, with two
  edges $e_a, e_b$ from a common vertex $v$. Then $G$ has a generic full-rank
  panel-hinge realization.
\end{lemma}
\begin{proof}
  \leanok
  \uses{lem:triangle-third-edge, lem:theorem-55-triangle, lem:triangle-normals,
        lem:genericity-device}
  \Cref{lem:triangle-third-edge} gives $V(G) = \{v,a,b\}$ and a completing edge
  $f : a$--$b$. \Cref{lem:triangle-normals} produces normals $n_0, n_1, n_2$ with
  pairwise nonvanishing joins and a linearly independent cyclic extensor family
  $C_0 = C(n_0,n_1)$, $C_1 = C(n_1,n_2)$, $C_2 = C(n_2,n_0)$. Assigning
  $v \mapsto n_0$, $a \mapsto n_1$, $b \mapsto n_2$ yields a seed $q_0$; each
  triangle edge's support extensor equals $\pm C_i$ (sign from the canonical
  orientation), so \cref{lem:theorem-55-triangle} gives rigidity on
  $\{v,a,b\} = V(G)$ at $q_0$. The genericity device
  (\cref{lem:genericity-device}) then re-realizes this rigidity at a
  general-position seed, giving the generic
  full-rank realization.
\end{proof}
\begin{fmlnote}
  This standalone triangle assembly is off the recursion's live path: the
  Case-III triangle floor (\cref{lem:case-III}) closes $|V| = 3$ through the
  general cycle family \cref{lem:cycle-realization} at $m = 3$, of which the
  triangle is exactly the smallest instance (a triangle \emph{is} the $3$-cycle,
  and KT state Lemma~5.4 at $3 \le |V| \le D$). It is kept here, like the $d = 3$
  candidate dispatch (\cref{lem:case-III-candidate-dispatch-d3}), because its
  direct three-normal seed is the accessible worked instance of the general
  cyclic construction.
\end{fmlnote}

\medskip\noindent\emph{The general-$d$ chain dispatch (KT \S6.4.2,
Lemma~6.13).} KT generalize the $d = 3$ Lemma~6.10 to arbitrary dimension
with the remark that ``the proof strategy is exactly the same as $d = 3$''
(p.~692). Two pieces of that strategy do lift verbatim: the redundant-row
mechanism of Claim~6.11 (\cref{lem:case-III-claim-6-11}) applies at the base
split exactly as it does at $d = 3$, and the discriminator's finish is
Lemma~2.1 at general grade (\cref{lem:case-III-claim612-extensor-span}), the
same span argument that closes the fixed three-candidate Claim~6.12 above,
now applied to $k+2$ panel normals instead of four points. What KT's remark compresses is the
bridge between the two: at $d = 3$ the redundancy is tested against three
\emph{fixed} candidates, while at general $d$ it must first be transported,
with a sign change, across the $d$ panels of a chain --- KT's eqs.\
(6.57)--(6.66) --- before the same span argument can select one. We spell
out both pieces in full below.

The setting: $G$ simple, minimal $0$-dof, no proper rigid subgraph,
$|V| \ge 4$; the dichotomy (\cref{lem:chain-data-extract}) yields a
length-$d$ chain $v_0 v_1 \cdots v_d$ (\cref{def:chain-data}) whose interior
vertices all have degree $2$ (the short-cycle disjunct is closed directly by
\cref{lem:cycle-realization}). All $d$ candidate frameworks are re-views of
a \emph{single} splitting-off: split off once at $v_1$ along $(v_0, v_2)$
with a fresh short-circuit edge $e_0$, giving the base graph
$G_1 = G_{v_1}^{v_0 v_2}$, which the all-$k$ induction hypothesis realizes as
a generic framework. The $i$-th candidate re-attaches $v_1$ with its new
hinge drawn on a line inside the panel of the $i$-th chain vertex; KT's
isomorphisms $\rho_i$ (eqs.\ (6.54)--(6.56), ``exactly the same framework'')
identify candidate $i$ with candidate $1$ by an index shift along the
chain, so there is no second split and no second application of the
induction hypothesis --- the formalization models the shift by a single
relabeling isomorphism transporting the base data to each chain panel in
turn (\cref{lem:splitOff-ofNormals-relabel},
\cref{lem:splitOff-rigidityRows-relabel}) rather than repeating the
isomorphisms $\rho_i$ one at a time.

\begin{lemma}[The shared redundancy and the chain discriminator; KT Lemma~6.13, eqs.\ (6.57)--(6.67)]
  \label{lem:case-III-chain-discriminator}
  \lean{CombinatorialRigidity.Molecular.PanelHingeFramework.chainData_fire_discriminator}
  \leanok
  \uses{def:chain-data, def:rank-hypothesis}
  Let $G$ be a simple minimal $0$-dof-graph with $|V(G)| \ge 3$ at general
  dimension ($1 \le k$, $D = \binom{k+2}{2}$), let $v_0 v_1 \cdots v_d$
  ($d \ge 2$) be chain data on $G$ (\cref{def:chain-data}), and suppose the
  all-$k$ induction hypothesis at every strictly smaller minimal
  $k'$-dof-graph ($k' \ge 0$ arbitrary): the graph is nonempty, it carries a
  general-position, link-recording realization when it is simple, and a
  bare full-rank realization (\cref{def:rank-hypothesis}) in all cases. If the split-off
  $G_1 = G_{v_1}^{v_0 v_2}$ has deficiency $0$ and carries a generic
  full-rank realization, then there is a nonzero functional
  $\rho_0 \in \bigl(\bigwedge^k K^{k+2}\bigr)^*$ with $\rho_0(C(v_0 v_2)) =
  0$, and a chain index $i$ such that $\rho_0$ does not annihilate the join
  of the $i$-th candidate panel's normal with some transversal vector.
\end{lemma}
\begin{proof}
  \leanok
  \uses{lem:case-III-claim-6-11, lem:case-III-nested-rank-lower,
        lem:case-III-claim612-extensor-span,
        lem:case-III-claim612-orthseq-vanish, lem:case-III-claim612,
        lem:case-III-claim612-line-in-panel-union}
  The realization is taken at a seed chosen as a common non-root of finitely
  many nonzero polynomials in the seed coordinates, all fixed before the
  seed: the rank polynomial of the base graph $G_1$, the eq.-(6.22) rank
  polynomial of the vertex-removal $G - v_1$
  (\cref{lem:case-III-nested-rank-lower}), a general-position polynomial,
  and the $(k+1) \times (k+1)$ determinant of the chain candidate panels'
  normal rows; over the infinite field $K$ their product has a non-root.
  At that seed the base graph attains full rank, so Claim~6.11
  (\cref{lem:case-III-claim-6-11}), applied once at the base pair
  $(v_0, v_2)$, converts the eq.-(6.22) bound into the shared redundancy
  $\rho_0 \ne 0$ with $\rho_0(C(v_0 v_2)) = 0$. The $d + 1 = k + 2$ panel
  normals of the chain's candidate vertices are linearly independent at the
  chosen seed (the determinant factor is nonvanishing there), so
  Lemma~2.1 at general grade (\cref{lem:case-III-claim612-extensor-span})
  makes their $\binom{k+2}{2} = D$ joins span $\bigwedge^k K^{k+2}$. A
  nonzero functional cannot annihilate every element of a spanning set
  (\cref{lem:case-III-claim612-orthseq-vanish}), so this is Claim~6.12 at
  general grade (\cref{lem:case-III-claim612}): $\rho_0$ fails to annihilate
  some witness join. Reading that join as a panel-meet via the general
  point-join $\leftrightarrow$ panel-meet duality
  (\cref{lem:case-III-claim612-line-in-panel-union}) fixes the matched
  candidate index $i$ and a transversal direction for its line.
\end{proof}
\begin{fmlnote}
  The Lean statement returns further data used only by
  \cref{lem:case-III-chain-dispatch}'s two realization branches: an
  endpoint selector recording every link of $G - v_1$, the chosen seed
  itself, $\rho_0$'s expression as a combination of the $D - 1$ hinge
  functionals at $(v_0, v_2)$, and an independent family of
  $D(|V(G_1)| - 1)$ rigidity rows of $G - v_1$ certifying the base rank ---
  Lean bookkeeping, not further mathematical content.
\end{fmlnote}

\begin{lemma}[The general-$d$ chain dispatch; KT Lemma~6.13]
  \label{lem:case-III-chain-dispatch}
  \lean{CombinatorialRigidity.Molecular.PanelHingeFramework.chainData_dispatch}
  \leanok
  \uses{lem:case-III-chain-discriminator, def:rank-hypothesis}
  With $G$, the chain data $v_0 v_1 \cdots v_d$, and the induction hypothesis
  as in \cref{lem:case-III-chain-discriminator}: if the split-off
  $G_1 = G_{v_1}^{v_0 v_2}$ has deficiency $0$ and carries a generic
  full-rank realization, then $G$ carries a generic full-rank realization.
\end{lemma}
\begin{proof}
  \leanok
  \uses{lem:case-III-chain-discriminator, lem:case-II-realization-placement,
        lem:splitOff-ofNormals-relabel, lem:splitOff-rigidityRows-relabel,
        lem:genericity-device}
  Apply \cref{lem:case-III-chain-discriminator} once, at the base split, to
  obtain the shared redundancy $\rho_0$ and the matched candidate index $i$.
  If $i$ is one of the two base-adjacent candidates, realize $G$ directly at
  the base split's own roles: $\rho_0$'s matched line closes the eq.-(6.29)
  top-left block of the eq.-(6.12) placement
  (\cref{lem:case-II-realization-placement}) to full rank. If $i$ is
  interior, transport $\rho_0$ to the $i$-th panel by the relabeling
  isomorphism identifying that panel with the base split
  (\cref{lem:splitOff-ofNormals-relabel},
  \cref{lem:splitOff-rigidityRows-relabel}): the degree-$2$ closure at each
  interior chain vertex (\cref{def:chain-data}) makes each step of the
  transport a column operation that reads $\rho_0$, with a sign change, at
  the next panel, until it reaches panel $i$. Either way the matched
  candidate attains rank $D(|V(G)| - 1)$ at a degenerate placement, and the
  genericity device (\cref{lem:genericity-device}) re-realizes that rank at
  a general-position seed.
\end{proof}

\begin{lemma}[Case III: $k = 0$, no proper rigid subgraph; KT Lemmas 6.10/6.13]
  \label{lem:case-III}
  \lean{CombinatorialRigidity.Molecular.PanelHingeFramework.case_III_realization_all_k,
        CombinatorialRigidity.Molecular.PanelHingeFramework.case_III_realization}
  \leanok
  \uses{def:panel-hinge-framework, def:rank-hypothesis,
        lem:simple-minimal-noRigid, lem:reduction-step}
  Let $G$ be a simple minimal $0$-dof-graph on $|V| \ge 3$ bodies with no
  proper rigid subgraph, at general dimension ($1 \le k$,
  $3 \le D = \binom{n+1}{2} = \binom{k+2}{2}$), and suppose the all-$k$
  induction hypothesis at every strictly smaller minimal $k'$-dof-graph
  ($k' \ge 0$ arbitrary): the graph is nonempty, it carries a
  general-position, link-recording realization when it is simple, and a
  bare full-rank realization (\cref{def:rank-hypothesis}) in all cases.
  Then $G$ carries a general-position full-rank realization at rank
  $D(|V|-1)$ --- the general-position conjunct of the zero-deficiency
  splitting-off case that the all-$k$ Theorem~5.5 induction needs. One
  hypothesis rides explicitly: a fresh edge
  name for $G$ itself (an unused label outside $E(G)$, needed for the
  splitting-off's short-circuit edge $e_0$). The eq.-(6.22) nested-IH rank
  bound is discharged internally from the all-$k$ IH
  (\cref{lem:case-III-nested-rank-lower}), not carried; the $d = 3$ form is
  the $k = 2$ specialization, and the floor $D \ge 3$ includes the planar
  case $n = 2$.
\end{lemma}
\begin{proof}
  \leanok
  \uses{lem:triangle-third-edge, lem:chain-data-extract,
        lem:cycle-realization, def:cycle-data,
        lem:simple-minimal-noRigid, lem:reduction-step,
        lem:case-III-chain-dispatch, lem:low-degree-vertex}
  Induction on $|V(G)|$, proving both conjuncts of the statement together.
  At $|V| = 3$ the recursion bottoms on the triangle floor. A degree-2 vertex $v$ exists by the
  average-degree count \cref{lem:low-degree-vertex}
  (\texttt{Graph.exists\_degree\_eq\_two\_of\_noRigid}); any other vertex $a \ne v$ of the
  three-vertex graph serves as its companion, since adjacency of $v, a$ is never used below (the
  adjacent-pair lemma \cref{lem:adjacent-degree-two-pair} is false at $D = 3$). The degree-2
  closure at $v$ supplies its two edges $e_a : v$--$a$, $e_b : v$--$b$, which pin the
  third edge and the vertex set $V(G) = \{v,a,b\}$ (\cref{lem:triangle-third-edge}); packaging
  $G$ as the $3$-cycle $C_3$ then realizes it as the $m = 3$ instance of the KT Lemma-5.4
  cycle family (\cref{lem:cycle-realization}, stated by KT at $3 \le |V| \le D$, of which the
  triangle is the smallest instance). At $|V| \ge 4$, the
  general-$d$ chain extractor (\cref{lem:chain-data-extract}) runs the KT
  Lemma-4.6/4.8 chain-vs-short-cycle dichotomy. On the short-cycle disjunct,
  the Lemma-5.4 cycle family (\cref{lem:cycle-realization}) closes $G$
  directly --- a genuine branch at $d \ge 4$, mathematically vacuous at
  $d \le 3$ (a cycle witness forces $|V(G)| = m \le d \le 3$, contradicting the
  $|V| \ge 4$ floor this disjunct runs under). On the chain disjunct it
  returns a length-$d$ chain $v_0 v_1 \cdots v_d$ together with the base
  $v_1$-split's minimality/simplicity/measure data; splitting off at $v_1$
  along $(v_0, v_2)$ gives a smaller simple minimal $0$-dof-graph
  (\cref{lem:reduction-step}; simplicity of the split-off holds at
  $|V| \ge 4$), so the induction hypothesis hands a generic full-rank
  realization of the base split $G_1$. The chain dispatch
  (\cref{lem:case-III-chain-dispatch}) lifts this to a generic full-rank
  realization of $G$ itself, establishing both conjuncts for $G$. At $d = 3$
  ($k = 2$) this chain dispatch is the fixed three-candidate case, worked out
  explicitly as an accessible entry point in
  \cref{lem:case-III-candidate-dispatch-d3}.

  The triangle base needs a single degree-2 vertex, not the $d = 3$-only
  \cref{lem:adjacent-degree-two-pair}. The dichotomy
  itself is run by the general extractor at every grade, its $|V| \ge 4$ cycle disjunct closed
  by \cref{lem:cycle-realization} (a genuine branch at $d \ge 4$, vacuous at $d \le 3$); and the
  cycle family is load-bearing for the $k > 0$ Case~II (\cref{lem:case-II}) as well.
\end{proof}
